\documentclass[11pt,a4paper]{article}


\usepackage[T1]{fontenc}
\setlength{\bigskipamount}{5ex plus1.5ex minus 2ex}
\setlength{\textheight}{24.45cm} 
\setlength{\textwidth}{15cm}
\setlength{\hoffset}{-1.2cm}
\setlength{\voffset}{-2.3cm} 
\usepackage{pdfsync} 
\usepackage[backend=biber]{biblatex}
\addbibresource{literature.bib}
\usepackage{amsfonts,amsmath,amsthm}
\setlength{\marginparwidth}{2cm}
\newtheorem{theorem}{Theorem}
\numberwithin{theorem}{section}
\newtheorem{lemma}[theorem]{Lemma}
\newtheorem{proposition}[theorem]{Proposition}
\newtheorem{conjecture}[theorem]{Conjecture}
\newtheorem{remark}[theorem]{Remark}
\usepackage{todonotes}
\theoremstyle{definition}
\newtheorem{definition}[theorem]{Definition}

\numberwithin{equation}{section}
\usepackage{amssymb}
\usepackage{hyperref}
\usepackage{graphicx}
\usepackage{wrapfig}
\usepackage{subfig}
\usepackage{epsfig}
\usepackage{float,mathrsfs}
\usepackage{bm}
\usepackage{enumerate}
\usepackage{enumitem}
\usepackage{bbm}
%
\newtheorem{algorithm}{Algorithm}
\usepackage{algpseudocode}
\usepackage[font=scriptsize]{caption}
\counterwithout{figure}{section}
\usepackage{color}
\definecolor{refkey}{rgb}{0.9451,0.2706,0.4941}\definecolor{labelkey}{rgb}{0.9451,0.2706,0.4941}

\definecolor{darkred}{RGB}{139,0,0}
\definecolor{darkgreen}{RGB}{0,100,0}
\definecolor{darkmagenta}{RGB}{139,0,139}
\definecolor{gray}{RGB}{180,180,180}


\newcommand{\setu}{{\mathrm{\mathfrak{u}}}}
\newcommand{\setv}{{\mathrm{\mathfrak{v}}}}
\newcommand{\setw}{{\mathrm{\mathfrak{w}}}}

\newcommand{\bsx}{{\boldsymbol{x}}}
\newcommand{\bsxi}{{\boldsymbol{\xi}}}
\newcommand{\bsu}{{\boldsymbol{u}}}

\newcommand{\bsy}{{\boldsymbol{y}}}
\newcommand{\brho}{{\boldsymbol{\rho}}}
\newcommand{\bsF}{{\boldsymbol{F}}}
\newcommand{\bsR}{{\boldsymbol{R}}}
\newcommand{\bsI}{{\boldsymbol{I}}}

\newcommand{\bsY}{{\boldsymbol{Y}}}
\newcommand{\bsC}{{\boldsymbol{C}}}
\newcommand{\bsgamma}{{\boldsymbol{\gamma}}}

\newcommand{\bsz}{{\boldsymbol{z}}}
\newcommand{\bst}{{\boldsymbol{t}}}
\newcommand{\bsg}{{\boldsymbol{g}}}
\newcommand{\bsomega}{{\boldsymbol{\omega}}}

\newcommand{\bsDelta}{\boldsymbol{\Delta}}
\newcommand{\rd}{{\mathrm{d}}}
\newcommand{\Dbf}{{\mathbf{D}}}
\newcommand{\bsphi}{\mathbf{\phi}}
\newcommand{\bspsi}{\mathbf{\psi}}

\newcommand{\Div}{{\mathrm{div}}}
\newcommand{\bbZ}{{\mathbb{Z}}}
\newcommand{\bbN}{{\mathbb{N}}}
\newcommand{\bfu}{\mathbf{u}}
\newcommand{\bO}{\mathbf{0}}
\newcommand{\dbsx}{\mathrm{d}\bfx}

\newcommand{\dxi}{\mathop{d\boldsymbol{\xi}}}
\newcommand{\dx}{\mathop{d\boldsymbol{x}}}

\newcommand{\bfv}{\mathbf{v}}

\newcommand{\bfy}{\mathbf{y}}
\newcommand{\bfx}{\mathbf{x}}
\newcommand{\bfg}{\mathbf{g}}
\newcommand{\bfS}{\mathbf{S}}
\newcommand{\bfxi}{\mathbf{\xi}}

\newcommand{\bfO}{\mathbf{0}}
\newcommand{\x}{\mathbf{x}}

\newcommand{\vxi}{\vec{\xi}}
\newcommand{\vg}{\vec{g}}
\newcommand{\va}{\vec{a}}
\newcommand{\vn}{\vec{n}}
\newcommand{\vbu}{\vec{\mathbf{u}}}
\newcommand{\vbv}{\vec{\mathbf{v}}}
\newcommand{\bsigma}{\mathbf{\sigma}}
\newcommand{\bssigma}{\boldsymbol{\sigma}}

\newcommand{\bSigma}{\mathbf{\Sigma}}

\newcommand{\uxi}{\underline{\xi}}

\newcommand{\bbR}{{\mathbb{R}}}
\newcommand{\satop}[2]{\stackrel{\scriptstyle{#1}}{\scriptstyle{#2}}}
\newcommand{\bstau}{{\boldsymbol{\tau}}}
\newcommand{\bsalpha}{{\boldsymbol{\alpha}}}
\newcommand{\bsbeta}{{\boldsymbol{\beta}}}
\newcommand{\bseta}{{\boldsymbol{\eta}}}
\newcommand{\bsnu}{{\boldsymbol{\nu}}}
\newcommand{\bsmu}{{\boldsymbol{\mu}}}
\newcommand{\bsOmega}{{\boldsymbol{\Omega}}}
\newcommand{\bskappa}{{\boldsymbol{\kappa}}}
\newcommand{\bszero}{{\boldsymbol{0}}}
\newcommand{\bsone}{{\boldsymbol{1}}}
\newcommand{\bsb}{{\boldsymbol{b}}}
\newcommand{\bsm}{{\boldsymbol{m}}}
\newcommand{\bsr}{{\boldsymbol{r}}}
\newcommand{\bsk}{{\boldsymbol{k}}}
\newcommand{\bsv}{{\boldsymbol{v}}}
\newcommand{\bsn}{{\boldsymbol{n}}}
\newcommand{\fallen}{ \left [\frac 1 2 \right] }
\newcommand{\fallennudelta}{ \fallfac{\delta}{\vnu}}

\newcommand{\bsp}{{\boldsymbol{p}}}
\newcommand{\bsq}{{\boldsymbol{q}}}
\newcommand{\bsw}{{\boldsymbol{w}}}
\newcommand{\bsell}{{\boldsymbol{\ell}}}
\newcommand{\bsS}{{\boldsymbol{S}}}
\newcommand{\bsW}{{\boldsymbol{W}}}
\newcommand{\bsE}{{\boldsymbol{E}}}
\newcommand{\ri}{\mathrm{i}}
\newcommand{\DivM}{\underline{\mathrm{div}}}
\newcommand{\D}{\partial}
\newcommand{\Dnuy}{\partial^{\bsnu}_{\bsy}}

\newcommand{\calU}{\mathcal{U}}
\newcommand{\calA}{\mathcal{A}}
\newcommand{\calM}{\mathcal{M}}
\newcommand{\calF}{\mathcal{F}}
\newcommand{\calT}{\mathcal{T}}

\newcommand{\N}{\mathbb{N}}
\newcommand{\R}{\mathbb{R}}

\newcommand{\bbF}{\mathbb{F}}
\newcommand{\bbE}{\mathbb{E}}
\newcommand{\calO}{\mathcal{O}}\newcommand{\calC}{\mathcal{C}}
\newcommand{\calD}{\mathcal{D}}
\newcommand{\calE}{\mathcal{E}}
\newcommand{\calB}{\mathcal{B}}

\newcommand{\calL}{\mathcal{L}}
\newcommand{\calH}{\mathcal{H}}
\newcommand{\calW}{\mathcal{W}}
\newcommand{\tr}{{\tt T}}
\newcommand{\esup}{\operatornamewithlimits{ess\,sup}}
\newcommand{\mix}{\mathrm{mix}}
\newcommand{\wor}{\mathrm{wor}}
\newcommand{\setD}{{\mathfrak{D}}}
\newcommand{\setS}{{\mathcal{S}}}
\newcommand{\mask}[1]{{}}
\newtheorem{assumption}[theorem]{Assumption}
\newtheorem{problem}[theorem]{Problem}
\newtheorem{example}[theorem]{Example}

\newcommand{\op}{\mathrm{op}}
\newcommand{\ess}{\mathrm{ess}}

\usepackage{pgfplots}
\pgfplotsset{compat=1.18}
\newcommand{\rev}[1]{{{#1}}}
\newcommand\DD[2]{%
\frac{\partial #1}{\partial #2}
}
\newcommand\DDt[2]{%
\frac{\partial^2 #1}{\partial #2}
}

\newcommand{\gdnota}[1]{\left[\tfrac{1}{2} \right]_{#1} }
\newcommand{\fallfac}[2]{\left[\tfrac{1}{2} \right]^{(#1)}_{#2} }
\newcommand{\fallfact}[1]{{\left[\tfrac{1}{2}\right]_{#1}}}
\newcommand{\norm}[2]{{\left\| #1 \right\|_{#2}}}

\newcommand{\cB}{{\mathcal{B}}}
\newcommand{\mC}{{\mathbb{C}}}

\newcommand{\ds}{\mathrm{ds}}
\newcommand{\brac}[1]{\left(#1\right)}
\newcommand{\abs}[1]{{\left|#1\right|}}
\newcommand{\NN}{{\mathbb N}}
\newcommand{\RR}{{\mathbb N}}
\newcommand{\vnu}{{\boldsymbol{\nu}}}
\newcommand{\veta}{{\boldsymbol{\eta}}}
\newcommand{\vm}{{\boldsymbol{m}}}
\newcommand{\vb}{{\boldsymbol{b}}}
\newcommand{\vzero}{{\boldsymbol{0}}}
\newcommand{\ve}{{\boldsymbol{e}}}
\newcommand{\cF}{{\mathcal F}}
\newcommand{\valpha}{{\boldsymbol{\alpha}}}
\newcommand{\vbeta}{\boldsymbol{\beta}}
\newcommand{\vy}{{\mathbf{y}}}
\newcommand{\bsdelta}{{\mathbf{\delta}}}

\newcommand{\vR}{{\boldsymbol{\rho}}}
\newcommand{\sett}[1]{\left\{#1\right\}}
\newcommand{\supp}{\operatorname{supp}}



\newcommand{\tododuc}[2][]{\todo[color=red!20,#1]{#2}}
\newcommand{\todoharri}[2][]{\todo[color=green!20,#1]{#2}}
\newcommand{\todotung}[2][]{\todo[color=blue!20,#1]{#2}}
\newcommand{\ac}[2][]{\todo[color=orange!40,#1]{#2}}
\title{Uncertainty Quantification for the Critical Buckling Load for Shells with Random Domain Perturbations}
%\title{Regularity of the Critical Buckling Load of Shells with Uncertain Domain}    

\author{Harri Hakula\footnotemark[2]\and Duc Khuat\footnotemark[2]\and Alexey Chernov\footnotemark[1]\and Tung Le\footnotemark[1] }


\renewcommand{\thefootnote}{\fnsymbol{footnote}}


\allowdisplaybreaks

\begin{document}
\maketitle

\begin{abstract} 

    Thin-walled structures are prone to large-magnitude deformations under compressive loadings. 
    The process where an unstable initial small displacement transitions into a large-scale 
    deformation is called buckling; the corresponding load at which this phenomenon occurs is the critical buckling load.
    %Despite the advances in engineering, manufacturing and mathematical modeling, manufacturing and 
    %design specifications from the 70s remain the most prominent source of recommendations\textemdash
    %even when 
    The prediction of the critical buckling load remains difficult due to the high 
    sensitivity to geometric imperfections. 
    %Despite its practical importance for lightweight engineering and manufacturing
    %the critical buckling load under random domain perturbations has not been analyed theoretically. 
    %this sensitivity has not been 
    %analyzed with the means of classical uncertainty quantification analysis.
   % Despite its importance, the discrepancy of theoretically predicted and observed buckling 
    %load of thin-walled structures prominently as large as 80\%.
    %It is well recognized that geometric imperfections are the primary reason for this discrepancy.
    Here, we provide a general framework to include random domain perturbations as a 
    random domain mapping and analyze the associated 
    general linear eigenvalue problem. The smallest eigenvalue of this system corresponds to the critical buckling load.
    We show that if the random domain mapping is Gevrey-regular and uniformly diffeomorphic, the
    corresponding, non-compact spectral operator remains Gevrey-regular to explicitly provided constants.
    Moreover, using the resolvent of the spectral operator, we show that if the eigenvalue cluster that 
    contains the smallest eigenvalue is isolated, Gevrey-regularity is preserved. Subsequently, 
    we establish a criterium for when this condition for the smallest eigenvalue cluster 
    is satisfied.
    This criterium depends on the ratio between total variation of the operator and 
    the initial spectral gap of the perfect, unperturbed domain.
    In the regime at which this criterium is not satisfied we observe numerically that
    eigenmodes with a high-wave number mix with the smallest eigenvalue cluster
    accompanied with a collapse of the critical buckling load. 
    It remains open to quantify this phenomenon.
   % Several challenges are addressed that are non-standard in the uncertainty quantification literature:
   % the spectral operator is not assumed to be compact, the eigenvalue is not isolated and we include general 
   % Gevrey-regular functions, particularly, random domain mappings are not assumed to be
   % analytically dependent on the parameter.


\end{abstract}
%

\footnotetext[2]{Department of Mathematics and 
Systems Analysis, Aalto University, FI-00076 Aalto, Finland ({\tt harri.hakula@aalto.fi}, {\tt duc.khuat@aalto.fi}).}
\footnotetext[1]{Institut für Mathematik Carl von 
Ossietzky Universität Oldenburg, Oldenburg, D-26111, Germany ({\tt alexey.chernov@uni-oldenburg.de }, {\tt tung.le@uol.de}).}


%\listoftodos
%\tableofcontents
\section{Introduction}
\paragraph{Motivation and Context.} 
The phenomenon of buckling of thin-walled structures (shell) describes the transition from 
an initial unstable configuration to a possibly large-magnitude
 displacement under small increase in loading. This so-called buckling response 
 often causes either a significant reduction in stiffness or a complete structural failure.
Relevant examples of shells that are prone to buckling are lightweight air- or 
spacecrafts, submarines and automotives \cite{WAGNER2020106451}. 
Notoriously, in a major series of laboratory experiments by NASA in the 1970s 
\cite{NASA_SP_8007_2020}, the knockdown factor, 
the ratio between observed and theoretical critical buckling load, was shown 
to be as low as twenty per cent. 
Despite the advances in computational resources, mathematical models and 
manufacturing technology, this report remains the most prominent source of 
recommendations for the design of buckling-critical shells. 

%And thus, one motivation is to improve such manufacturing specifications 
%by performing forward and inverse uncertainty quantification for the critical buckling load.
%sharpening 
%the bound on the critical bucklo through a general treatment of domain uncertainty
%propagated to the critical buckling load.
%and extending the quantification of sensitivity to 
%different geometries and geometric imperfections. 
It is well recognized that geometric imperfections are the primary 
reason for the large discrepancy between the observed and 
theoretical critical buckling load \cite{WAGNER2020106451}. 
Naturally, efforts have been directed to the theoretical and 
numerical analysis of geometric imperfection sensitivity of the critical buckling model 
\cite{HUTCHINSON2018157,Harutyunyan2023,HILBURGER2004369, SCHENK20073424}. 
In this article, we apply a general uncertainty quantification analysis of 
the generalized eigenvalue problem that results from the linear buckling model of general shell structures. 
Moreover, we show Gevrey-regularity of the corresponding 
critical buckling load under the assumption that the random domain perturbations are 
Gevrey-regular and uniformly diffeomorphic. The accompanied 
constants are central to inform higher-order convergence methods 
such as quasi Monte-Carlo or 
sparse grid approximation methods \cite{kuo2016applicationquasimontecarlomethods,BieriDiss, Haji_Ali_2016}. 
Consequently, this result enables standard 
applications of methods in forward uncertainty quantification, 
inverse problems, and topological and structural as well as 
experimental design optimization \cite{kaarnioja2026quasimontecarlobayesiandesign,LatifiRostami2019}.


\paragraph{Problem statement.} 
The eigenvalue problem associated to the buckling problem is obtained from considering a 
linearization of the second variation of the elastic energy functional. 
With an increase in load, initially stable configurations may 
change sign in the second variation of energy and thus become unstable. 
Subsequently, a structure in an unstable configuration can suddenly 
transition to a large-magnitude deformation. In this linearization 
the pre-stress components involved in the non-linear term of the second 
variation of elastic energy depend linearly on the traction 
load and can hence be factored out as a dimensionless constant. The 
resulting equations resemble a general eigenvalue problem where 
the lowest eigenvalue corresponds to the dimensionless critical buckling load. 
Here, the considered model is based on a geometric non-linear 
deformation energy function and was analyzed, for example, in 
\cite{DaugeandSuriNumericalapproximation, SZABO199943}.

The parametrized general eigenvalue problem reads as:
\emph{Find the smallest value $\lambda(\bsy) $ and $\bsu(\bsy)$ such that}
\begin{align}\label{eq:bucklingPreview}
    a_{\bsy}(\bsu(\bsy), \bsv) = \lambda(\bsy) b_{\bsy}(\bsu, \bsv) , \quad  \forall \bsv \in V,
\end{align}
where $a_\bsy$ is the variational bilinear form for linear elasticity, $b_\bsy$ is the pre-stress bilinear form and 
$\bsy \in \calU$ is a parameter for the random domain mapping that models domain uncertainty in the shell. 
Perturbations in the domain are modeled with the framework of random domain mappings \cite{Harbrecht2023DomainUncertainty} 
and enter the equation through a pull-back to the nominal domain.
Since parametrized eigenvalues are known to admit non-smooth dependencies, 
we need to restrict our analysis on eigenvalue clusters. 
Eigenvalue cluster are subsets of eigenvalues that remain isolated; that is, 
the eigenvalues of the corresponding eigenmodes do 
not cross with the eigenvalues of eigenmodes outside of the cluster for the whole parameter domain. 
We are concerned with the question of whether Gevrey-regularity of the 
random domain is preserved for the eigenvalue cluster, 
in the sense that its spectral operator is Gevrey-regular. 
Secondly, we want to quantify how large the eigenvalue cluster is 
and if it allows for an approximation of the critical buckling load.
 
\paragraph{State of the art and the Research gap.}
Parametrized eigenvalue problems are well studied, see for example 
\cite{grubišić2021stochasticcollocationmethodcomputing,chernov2024analyticgevreyclassregularity} 
and the references therein. However, most existing results are only valid for compact operators for which the whole spectrum 
is discrete and where the eigenvalue of interest is isolated. The buckling problem, however, admits a few challenges: 
With the Riesz representation theorem on \eqref{eq:bucklingPreview}, we obtain the linear operator general 
eigenvalue problem
\begin{align*}
    \calA_\bfy \bfu &= \lambda(\bfy)\calM_{\bfy} \bfu,
\end{align*}
where $\calA_\bsy$ is uniformly invertible and $\calM_\bsy$ is possibly indefinite.  
Moreover, since $a_\bsy$ and $b_\bsy$ are a vector-valued first order bilinear forms, 
the spectral operator $\calT_\bsy = \calA_\bsy^{-1} \calM_\bsy$ is non-compact. It is known by 
\cite{DaugeandSuriNumericalapproximation} that the non-compact spectral operator from \eqref{eq:bucklingPreview} 
is prone to spectral pollution; whole intervals can become eigenvalues and thus numerical approximations 
seem to not converge in a certain range of eigenvalues. In \cite{DaugeandSuriNumericalapproximation}
has shown that the smallest eigenvalue can be approximated pollution-free whenever the structure 
is sufficiently thin. 

Previous work on the buckling phenomenon has focused on the propagation of the effect of small idealized perturbations, 
based on analytic extensions of the dimension reduced model solution, on the critical buckling 
load in  \cite{HUTCHINSON2018157}. In related work \cite{Harutyunyan2023}, 
the authors proved that the observed asymptotic 
behavior of critical buckling load changes 
under curvature variations and that the critical load decays 
with a larger rate than the theoretic rate for a perfect geometry 
when the thickness tends to zero. % \cite{dauge:hal-00012083}.
On the numerical side, in \cite{SCHENK20073424} the authors have conducted extensive 
numerical experiments with thin-walled cylinders with geometric imperfections that matches the manufacturing specimens
made available by the initial imperfection data bank in \cite{arbocz1979initial}. It is shown in \cite{HILBURGER2004369} 
that the prediction gap can be as low as 5\% percent when geometric imperfections are included.
To the best of our knowledge, general geometric imperfections have not been considered 
theoretically in the context of buckling.
%Here, we analyze the propagation of random domain perturbations through the associated linear eigenvalue
 %problem using the general uncertainty quantification methodological framework.

\paragraph{Contribution.}
 %We analyze the propagation of random domain perturbations through the associated linear eigenvalue
%problem using the general uncertainty quantification methodological framework. 
Our main result is that the
eigenvalue cluster of the smallest eigenvalue, that is associated to the critical buckling load, 
preserves Gevrey-regularity with an explicitly provided constant.
%when the domain perturbation function is Gevrey-regular 
%and its total variation is bounded 
We introduce several new proof techniques, 
further leveraging the alternating factorial 
technique \cite{chernov2024analyticgevreyclassregularity}. 
We show Gevrey-regularity of the spectral operator by differentiating
its resolvent \cite{kato1995perturbation}. 
%The resolvent method can be used for the 
%computation of the eigenspace and further exploited for the spectral analysis.

Another open question is the asymptotic analysis 
for the critical buckling load, that is, the 
behaviour of the spectrum when the thickness of 
the structure goes to zero. In this limit case, 
we observe in numerical simulation that mixing 
of eigenmodes with high wave numbers with the 
smallest eigenvalue occurs. It remains outside 
of the scope of this work to combine the 
asymptotic analysis with the parametrized 
general eigenvalue problem.






\paragraph{Outline.}
The linear buckling model with domain perturbation is described in Section \ref{sec:LinearModel}. 
In Section \ref{section:AbstractResults}, we prove abstract results on the 
spectrum of $\calT_{\bfy}$ and particularly, prove that an isolated 
eigenvalue cluster of $\calT_\bfy$ preserves Gevrey-regularity. %provided that  $\bfy \mapsto \calT_\bfy$ is Gevrey-regular. 
In Section \ref{sec:regularityOfTy}, we show that $\calT_\bfy$ is 
Gevrey-regular if the random domain mapping is Gevrey-regular and uniformly 
diffeomorphic. Morevover, we provide the constants of the regularity statement. 
We conclude with numerical experiments in Section 
\ref{sec:numerical-example} and a discussion in Section \ref{sec:Discussion}.




 
\section{The Linear Model for Buckling with Random Domain Perturbations}\label{sec:LinearModel}
\subsection{Random Domain Perturbations}
Let $\Omega \in \R^3$ be a reference domain and let $(\mathcal{R},\calA,P)$ be a probability space. 

A common approach  to model random domain perturbations \cite{Harbrecht2023DomainUncertainty}
is to associate a random domain $D[\bsomega]$ with random parameters $\boldsymbol{Y}: \mathcal{R} \to \calU, \omega \mapsto \boldsymbol{Y}[\bsomega]$, where $\calU = [-1/2,1/2]^{\N^*}$, and the mapping  
\begin{align}\label{eq:definitionPhiandJ}
    \Phi: \calU \to C^2(\Omega;\R^3), \quad J(\bsy) := \DD{\Phi}{\bsxi}(\bsy)= \begin{pmatrix}
        \DD{\Phi}{\xi_1}& \DD{\Phi}{\xi_2}& \DD{\Phi}{\xi_3}
    \end{pmatrix} (\bsy)
\end{align} such that $\Phi$ uniformly twice-continuous differentiable and $C^1$-diffeomorphism, so that for all $\bsy \in \calU$
\begin{align*}
    \| \Phi(\bfy) \|_{C^1(\Omega;\R^3 )} \leq c_{\Phi} \quad \text{and} \quad  \| \Phi^{-1}(\bfy) \|_{C^1(\Omega;\R^3 )} \leq c_{\Phi},
\end{align*} for some constant $c_\Phi \geq 1$. It is common to suppress the dependence on $\bsomega$ and write $ D(\bsy)  = D(\boldsymbol{Y}[\bsomega])=  D[\bsomega]$ for $\bsy \in \calU$.
In practice, the random parameters $\boldsymbol{Y}[\bsomega]$ induce weights $\bsgamma \in \ell^2(\N)$ that are $\ell^2$-summable. In fact, the weights often correspond to a Karhunen-Loève decomposition of a random field in its principal components. Note that, any random field with finite variance admits such a decomposition, however, the components of $\boldsymbol{Y}[\bsomega]$ are not necessarily bounded. This property is a common but not restrictive simplification, see, for example, the textbook by G. Lord and C. Powell \cite[Chapter 7]{Lord_Powell_Shardlow_2014}. Fixing such weights $\bsgamma \in \ell^2(\N)$ associated with the random domain mapping, we define the norm on $\calU $ as 
\begin{align*}
    \| \bsy \|_\calU := \sup_{i \in \N^*}  \boldsymbol{\gamma}_i \vert \bsy_i \vert.
\end{align*} 


\subsection{Linear Elasticity}\label{subsec:linearelasicity}
%Standard practice for the stability analysis of shell structures is the use of tensor notation and Einstein's summation convention. 
In the following, we describe the relevant notational conventions.
We denote the set of all tensors over $\R^3$ of order $d$ as $\mathbb{T}_d$. Order two tensors are called matrices.% based on the textbook by Chapelle and Bathe \cite[Chapter 2.1]{ChapelleBathe2011}. 
%\emph{Latin} indices $i,j,k,l,\ldots$ range from 1 to 3 and \emph{Greek} indices $\alpha,\beta,\gamma,\ldots$ range from 1 to 2 if not otherwise specified. 
%For $d \in \N$ tensors of order $d $ are denoted as $^{(d)}S $ with components $S^{i_1 \ldots i_d}$ in contravariant and $S_{i_1 \ldots i_d}$ in covariant form, respectively.
%We write first order tensors  $\bsu = \{ u_i \}$ in bold.
 %In Einstein's summation convention, we omit summation signs for all indices that appear once as a subscript and
%once as a superscript in an expression. Finally, we describe the set of all tensors over $\R^3$ of order $d$ as $\mathbb{T}_d$.
%For a basis $\bsg^{i}$, the tensor product $\bsg^i \otimes  \bsg^{j}$ is a basis for the second order tensors.
%For a fixed basis, a second order tensor is associated with entries $T_{ij}$ and is represented as $\boldsymbol{T}= T_{ij} \bsg^i \otimes  \bsg^{j}$. Higher order tensors are treated in a similar fashion.
%The dot product for higher order tensors is defined as taking the last order of the first argument and the first order of the second argument in a sum of products. That is,  for two second order tensors $\boldsymbol{T}$ and $\boldsymbol{M}$ defined over a basis $\{ \bsg^i \}$, the dot product is a  tensor of order two  
%\begin{align*}
%\boldsymbol{T} \cdot \boldsymbol{M} = T_{ij}M_{kl} (\bsg^j \cdot \bsg^k )\bsg^{i} \otimes \bsg^l.    
%\end{align*}

%The tensors associated to the domain transformation are the covariant and contravariant basis, the second order metric tensor $\bsg$ and the volume elements of $\Phi$. 
%In the following, all expressions that contain $\Phi$ are functions of $\bsy$.
%and for notational brevity, we omit this dependence. 
Let $\bsxi  = (\xi_1,\ldots,\xi_3) $ denote the coordinates in the reference domain  $\Omega$ and $\bsx = (x_1,x_2,x_3)$ be the coordinates in the physical domain  $ D(\bsy) = \Phi(\Omega,\bsy)$ for $\bsy \in \calU$. Let $\bsu,\bsv: \Omega \to \R^3$ be two vector-valued functions and denote $\hat \bsu = \bsu \circ \Phi^{-1}$.
With the notation
\begin{align*}
    \DD{\bsu}{\bsxi} =    \begin{pmatrix}
         \DD{u_1}{x_1}& \DD{u_1}{x_2}& \DD{u_1}{x_3}  \\    
         \DD{ u_2}{x_1}& \DD{ u_2}{x_2}& \DD{u_2}{x_3}\\  
         \DD{ u_3}{x_1}& \DD{u_3}{x_2}& \DD{ u_3}{x_3}
    \end{pmatrix}, \quad \nabla_{\bsxi} \bsu = \DD{ \bsu}{\bsxi} ^{\top}, \quad \text{for } i,j \in \{1,2,3\},
\end{align*}
and the chain rule, we have
 \begin{align}\label{eq:partialderivativesinX}
\DD{\hat \bsu}{\bsx} =  \DD{\bsu}{\bsxi} \DD{\Phi^{-1}}{\bsx} =  \DD{\hat \bsu}{\bsxi} J^{-1} \quad \text{and } \quad  \nabla_{\bsx} \hat \bsu = J^{-\top} \nabla_{\bsxi} \bsu.
\end{align} Thus the gradient with respect to $\bsx$ is dependent on $\bsy \in \calU$.

To define the weak formulation of linear elasticity with random domain perturbations, we define the admissible energy space 
\begin{align}\label{eq:def_Vspace}
    V = \{ \bsu \in H^1(\Omega;\R^3) :  \bsu =  0  \text{ on }\partial \Omega_D   \}
\end{align}
equipped with the Sobolev $H^1$-norm, where we assume that we have a partition of the boundary $\partial \Omega  = \partial \Omega_D \sqcup \partial \Omega_T$ in non-zero area boundary parts where homogeneous Dirichlet boundary conditions and traction are specified. 
The linear strain is defined as 
\begin{align*}
    e_{\bsy}(\bsu) = \frac 1 2 \left(  \nabla_{\bsx} \hat \bsu +  (\nabla_{\bsx} \hat \bsu)^\top \right).
\end{align*}
Note that the gradient with respect to $\bsx$ depends on $\bsy \in \calU$ and hence does the linear strain on $D(\bsy)$.
For a fourth order tensor $\calC = \{ \calC_{ijkl}\}$ and $A,B \in \mathbb{T}_2$, we define the double dot product and the tensor-matrix product
\begin{align*}
    A : B &= \sum_{i,j=1}^3 A_{ij} B_{ij},  \\
    \calC : A :B &= \sum_{i,j,k,l=1}^3 \calC_{ijkl} A_{ij}B_{kl}, \\
    (A\calC)_{ijkl} &= \sum_{r=1}^3 A_{ir} C_{rjkl}, \quad i,j,k,l\in \{1,2,3\}\\
    (\calC B)_{ijkl} &= \sum_{s=1}^3 B_{sl} C_{ijsl}, \quad i,j,k,l\in \{1,2,3\}.
\end{align*} 
Note that 
\begin{align*}
     A : (MBN) &= (M^\top A N^\top) : B , \\
     \calC : MA : NB &=  (A^\top \calC N) : A : B, 
\end{align*}
for all $A,B,M,N \in \mathbb{T}_2,$ and $ \calC \in \mathbb{T}_4.$
 The bilinear form for linear elasticity is defined on $V$ as
 %For a tensor valued function $\bsu(\bsxi)$, we associated the 
 %We define functions in tensor notation $\bsu = \{ u_i \}$ and denote the bilinear form for \emph{linear elasticity} as
 \begin{align*}
    a_\bsy(\bsu,\bsv) &:=   \int_{D(\bsy)} \hat{\calC}:  e_\bsy( \hat \bsu) : e_\bsy( \hat \bsv)  \dx, \quad \bsy \in \calU,
\end{align*} where  $\hat {\calC}= \calC \circ \Phi $ and $ \calC= \{ \calC_{ijkl} \}$ is Hooke's material stiffness tensor defined as functions on the reference domain $\Omega$ and is symmetric in the sense of
\begin{align}\label{eq:symmetryC}
    \calC_{ijkl} = \calC_{jikl} = \calC_{ijlk} = \calC_{klij},\quad \text{ for all } i,j,k,l \in \{1,2,3\}.
\end{align}

Moreover, we assume that $\calC$ is a uniformly coercive and bounded stiffness tensor.
  That is, there is a constant $C_{\calC} > 0$ such that for all $\bfxi\in\Omega$,
  \begin{align}
  \calC(\bsxi): \boldsymbol{\alpha} : \boldsymbol{\alpha}& \geq C_{\calC}^{-1} \| \boldsymbol{\alpha} \|_F^2  &&\text{for all \emph{symmetric} matrices } \boldsymbol{\alpha} , \label{eq:stiffnesstensorcoercivity}\\
  \calC(\bsxi): \boldsymbol{\alpha} : \boldsymbol{\beta}& \geq C_{\calC} \| \boldsymbol{\alpha} \|_F  \| \boldsymbol{\beta} \|_F&&\text{for all matrices } \boldsymbol{\alpha},   \boldsymbol{\beta}\label{eq:stiffnesstensorupperbound}\\
  \end{align}
  where $\|\cdot\|_F$ is the Frobenius norm.
  
%To emphasize the dependence of the expression $\partial \hat v/ \partial x_l$ on $\bsy$, we may only consider the pulled back version using \eqref{eq:partialderivativesinX}
%\begin{align}\label{eq:linearelasticity_a_y}
%     a_\bsy(\bsu,\bsv) &:=   \int_{\Omega} \calC^{ijkl} \sqrt{g(\bfy)} \left(g^r_j(\bsy) u_{i,r}\right)  \left(g^s_l(\bsy) v_{k,s}\right) \dxi, \\ 
%     &= \int_{\Omega} \tilde{\calC}^{ijkl}(\bfy)  u_{i,j} v_{k,l}\dxi, 
%\end{align}
%where after exchanging indices $r$ and $j$, and $s$ and $l$, we have
%\begin{align*}
%    \tilde{\calC}^{ijkl}(\bsy) = \calC^{irks} g^j_r(\bsy) g^l_s(\bsy)  \sqrt{g(\bsy)} \in L^1(\Omega;\mathbb{T}_4), \quad \bsy \in \calU,
%\end{align*} where $\bsg^i$ and $g$ are mappings $\calU\to C^1(\Omega;\R^3)$.
%
After transforming the bilinear form to the reference domain and using \eqref{eq:partialderivativesinX}, we obtain a vector-valued stochastic diffusion bilinear form. The regularity of the solution of the corresponding variational formulation is well understood \cite{Reference}. The linear buckling problem, however, results from a coupled vector-valued stochastic diffusion equation and requires non-standard techniques. This coupling is described in the next Section \ref{section:linearbuckingformulation}. %€In Section \ref{section:AbstractResults} such a method is presented.
%, as we will see in the following Section.
%and requires non-standard techniques. 
%For example, a special case has been proven in \cite{Harbrecht2023DomainUncertainty}.

%and moreover satisfies
% \begin{align*}
%     C_{ijkl} \alpha_{ij} \alpha_{kl} \geq 0
% \end{align*}
% for any matrix $\{ \alpha_{ij} \}$.
 % \begin{align}
     %\| \bsu \|_V  := \left( \sum_{i,j = 1}^3  \| u_{i,j} \| ^2_{L_2(\Omega) } \right)
%\end{align}
\subsection{Linear Buckling Model with Random Coefficients}\label{section:linearbuckingformulation}
%We consider the three-dimensional linear buckling problem based on the potential energy criterion used in \cite{DaugeandSuriNumericalapproximation}.
Buckling in structural mechanics can be described as a \emph{bifurcation} phenomenon: as the traction load increases, the reference (prebuckling) solution of the governing equations loses stability, and the structure may switch to a different equilibrium branch. The bifurcation point of the traction load is the critical value at which the second variation of elastic energy at the reference solution branch ceases to be positive definite. Consequently, the \emph{linear buckling model} is obtained by retaining only the leading-order terms of this second variation equation. Here, we are concerned with the analysis of the linear buckling model based on the \cite{DaugeandSuriNumericalapproximation}; a detailed derivation applying the procedure described in this paragraph is given in \cite[Section~2]{Grabovsky2007TheFS} and \cite[Section 2]{HarutyunyanBuckling}.\tododuc{brief mentioning of limitations} After restating the problem in the physical domain $D(\bsy)$, we derive the random coefficient model on the reference domain $\Omega$.
%In the remainder of this section, we describe the linear bucking problem on $D(\bsy)$ and then restate %the problem on the reference domain $\Omega$.

%The random domain mapping formulation into a random parameter dependent linear eigenvalue problem  with random coefficients. % in the sense of \eqref{eq:linearelasticity_a_y}. 

%The mathematical formulation for the buckling \ac{buckling} problem arises from setting the second variation of elastic energy to %zero and.
% from considering only leading order terms in the equation where the second variation of elastic energy is set to zero. Mathematically, buckling occurs when the second derivative of the elastic energy function changes it sign; an initially stable solution becomes unstable and a bifurcation occurs. Consequently, the structure snaps to a new state. The goal of the stability analysis is to find the critical buckling load $\lambda_{\text{cr}}$,  at which this bifurcation occurs.
%For a derivation of the linear and non-linear buckling model, see, for example \cite[Section 2]{Grabovsky2007TheFS}. Therein, also limitations of the model \ac{model?} are discussed.

%\paragraph{Random Linear Buckling Problem.}
%The random domain linear buckling problem is retained by a classical coordinate transformation with the random domain mapping $\Phi$. 

The linearization of equations of non-linear elasticity requires a pre-stress weighting. This pre-stress is associated to the prebuckled stress state and is derived as follows:
%The pre-stress is derived from a random domain source problem:
%First, the reference solution is solved to obtain the pre-stress $\bssigma_*$ and then the linear eigenvalue problem can be formulation 
%First, we solve the pre-stress from the source problem where we assume negligible body load, that is, we assume that gravity is small compared to the traction load.
%Solving for the reference solution on the randomly perturbed domain, gives the following problem:
\emph{Find the pre-existing stress $\bssigma_0: \calU \to H^1(\Omega;\mathbb{T}_2)$ such that} 
\begin{align}\label{eq:sourceproblem_weak}
    \int_{D(\bsy)} \bssigma_0 : e_\bsy(\bsv) \dx  =\int_{D(\bsy)}  F_0\cdot  \bsv \dx  + \oint_{\partial D(\bsy)_T} T_0 \cdot  \bsv \mathop{dS}
\end{align}
for all $\bsv \in V$, where $\partial D(\bsy)_T = \Phi(\partial \Omega_T; \bsy)$, and body load $F_0$ and traction force are $T_0$ defined on the union of all $D(\bsy)$. 
As the body load is neglibile in our setting, we assume that $F_0$ is zero. The solution $\bssigma_0$ is then linear with respect to the traction load $T_0$, such that, 
for a dimensionless constant $\lambda \in \R$, we may write
\begin{align*}
    \bssigma_0= \lambda \bssigma_*, \quad T_0 =\lambda T
\end{align*} and solve for $\sigma_*$ on a unit traction load 
\begin{align}\label{eq:sourceproblem_weak_star}
    \int_{D(\bsy)} \bssigma_*: e_\bsy(\bsv) \dx  =  \oint_{\partial D(\bsy)_T} T \cdot  \bsv \mathop{dS}  =: \calF_{\bsy}(\bsv) .
\end{align}
Moreover, we assume that $\calF_{\bsy}$ is uniformly bounded, that is, there exists $C_{T} \geq 0$ such that for all $\bsy \in \calU$
\begin{align}\label{eq:boundOnTractionLoad}
    \calF_{\bsy}(\bsv) \leq C_T \| \bsv \|_{V}, \quad \forall v \in V.
\end{align}
The geometric stiffness bilinear form is defined by the energy of the weighted second order term of the nonlinear Green Lagrange strain tensor, as
\begin{align*}
    b_\bsy(\bsu,\bsv) &:= -\int_{D(\bsy)} \bssigma_*: (\nabla_{\bsx}\bsu)^\top  \nabla_{\bsx}\bsv \dx.
\end{align*} 
The parametrized linear buckling condition as in \cite{SZABO199943,DaugeandSuriNumericalapproximation,Grabovsky2007TheFS} is:
\begin{problem}\label{problem:linearbuckling}
   Find the smallest eigenvalue $\lambda_{\text{cr}}(\bsy)$ such that there exists $\bsu(\bsy) \in V$  
\begin{align}\label{eq:eigenvaluebucklingproblem}
    a_\bsy(\bsu(\bsy), \bsv )  = \lambda_{\text{cr}}(\bsy) b_\bsy(\bsu(\bsy), \bsv ), \quad \text{for all } \bsv \in V, \text{ and } \bsy \in \calU.
   \end{align}
\end{problem}
%For the random domain transformation, we adopt the notation by Harbrecht et al. in \cite{Harbrecht2023DomainUncertainty} to tensor notation. \ac{I'd better refer to a standard textbook, like Bathe}
% \ac{$\partial_j u_i$ is obviously not $\DD{u_i}{x_j}$. So what is this?}
%Let us define the Jacobian of the random domain mapping and its contravariant basis $J =  \{\Phi_{i,j} \}, J^{-\top} =  \{g^j_{i} \}$, where $g^j$ \ac{Introduce the metric tensor and use the notation consistently. Pay attention to covariant and contravariant coordinates (rising of an index, see Bathe)} is the contravariant basis.
%Using the identity
%\begin{align}\label{eq:partialderivativeuinU}
%   \partial^{\bsx}_j u_i = g^{r}_j u_{i,r},
%\end{align}
%the bilinear forms can be transformed back to the reference domain $\Omega$.
By using \eqref{eq:partialderivativesinX}, the symmetry of $\calC$ in the sense of \eqref{eq:symmetryC} and of $\bssigma_*$, and the transformation on the reference domain, we obtain the vector-valued random coefficient eigenvalue problem. 
%Note that both bilinear forms involve first-order derivatives. The corresponding spectral operator is therefore non-compact. 
\begin{problem}[Random Coefficient Linear Buckling Problem]\label{problem:transformed_linearbuckling}
    Find the smallest eigenpair $\lambda_{\text{cr}}(\bsy)$ and $\bsu(\bsy) \in V$ such that \eqref{eq:eigenvaluebucklingproblem} is satisfied, where the bilinear forms are
    \begin{align}\label{eq:bilinearforms}
    a_\bsy(\bsu, \bsv) &=  \int_{\Omega} \tilde{\calC} : \nabla_{\bsxi} \bsu :\nabla_{\bsxi} \bsv \dxi,  \\
    \label{eq:bilinearforms2}
    b_\bsy(\bsu, \bsv)  & = -\int_{\Omega}  \tilde{\bssigma}_* : (\nabla_{\bsxi} \bsu)^\top   \nabla_{\bsxi} \bsv  \dxi,
\end{align}
where \ac{Inconsistent, see Bathe}
    \begin{align}\label{eq:transformed_stiffness}
    \tilde \calC(\bsy)  &= \det J(\bsy) \left [J^{-1}(\bsy)\calC J^{-\top}(\bsy) \right] \quad \in L^1(\Omega;\mathbb{T}_4),  \\
    \label{eq:transforme_prestress}
    \tilde{\bssigma}_*(\bsy) &= \det J(\bsy)  \left[ J^{-1}(\bsy) \bssigma_*(\bsy) J^{-\top}(\bsy) \right] \quad \in H^1(\Omega;\mathbb{T}_2).
    \end{align}
 
%and where $J$ and $\bfg^{k}$ and $\sigma_*$ depend on $\bsy \in \calU$.
The unit pre-stress is obtained as the solution to \eqref{eq:sourceproblem_weak_star} after transforming the equation to $\Omega$. \emph{Find $\bssigma_* : \calU \to H^1(\Omega;\mathbb{T}_2)$ such that} 

\begin{align}\label{eq:sourceproblem_weak_pulled_back}
    \int_{\Omega} \tilde \bssigma_* : \nabla_{\bsxi} \bsv \dxi  =  \oint_{\partial \Omega_T} \tilde T : \bsv \mathop{dS},
\end{align}
for all $\bsv \in V$ and all $\bsy \in \calU$, where
\begin{align}
    \tilde \bssigma_*(\bsy) &=  \det J(\bsy) \left[ J^{-1}(\bsy) \bssigma_{*}(\bsy)  \right] \quad \quad\quad \in H^1(\Omega;\mathbb{T}_2),  \\
    \tilde T(\bsy) &=  \det J(\bsy)   \|J^{-\top}(\bsy) n^{\text{ref}}_r \|_2  T\circ \Phi(\bsy) \quad \in L^1(\Omega;\mathbb{T}_1),\label{eq:transformedtractionload}
\end{align}
where $\|  \cdot \|_2$ denotes standard Euclidean-norm on $\R^3$ and $\mathbf{n}^{\text{ref}} \in \R^3$ denotes the normal vector on $\partial \Omega$.
\end{problem}

%The dependence of $\sigma_*$ is defined by the source problem for the pre-stress. 
\begin{remark}[Non-compact spectral operators] \label{remark:noncompactOperator}
The linear random eigenvalue problem in Problem \ref{problem:transformed_linearbuckling} is a non-compact linear operator, since the bilinear forms $a_\bsy$ and $b_\bsy$ are of the same order.
Such operators pose problems in the numerical computation of the spectrum, since whole intervals become possible limits of eigenvalues in the finite dimensional approximation, polluting the spectrum and misinforming the interpretation of numerical results. 
However, the main result in \cite[Theorem 5.1]{DaugeandSuriNumericalapproximation} states an explicit criteria, such that a pollution-free eigenvalue approximation for the smallest eigenvalues on thin structures \ac{for the smallest eigenvalue?} is guaranteed. 
Consequently, our analysis of the critical eigenvalue that solves \eqref{eq:eigenvaluebucklingproblem} is well-posed, whenever the underlying domain is sufficiently thin.
\end{remark}

%Before we show the Gevrey-regularity of the bilinear forms $a_{\bsy}$ and $b_{\bsy}$, we describe 
\section{Abstract Results for Gevrey Regularity of Spectral Operators}\label{section:AbstractResults}
\tododuc{Write a more informative paragraph, connecting this better to the rest.}
%The parametrized eigenvalue problem in \eqref{eq:eigenvaluebucklingproblem} is a well studied problem.  
%We outline the analysis of the spectrum of \eqref{eq:eigenvaluebucklingproblem} with the use of the residual technique.
We describe how to quantify Gevrey-regularity of parametrized eigenvalue problems of the form in  \eqref{eq:eigenvaluebucklingproblem}. The following results and techniques are mostly well-established.
However, the proofs of the Gevrey-regularity of the spectral resolvent operators are novel contributions.
Here, we assume that the parametrization of $a_\bsy$ and $b_\bsy$ are Gevrey-regular and we provide a proof for this statement in Section \ref{sec:regularityOfTy}.
\subsection{Linear Eigenvalue Problem For Non-Compact Operators}\label{subsection:Tydefinition}
The setting is similar to \cite[Section 3]{DaugeandSuriNumericalapproximation}.
Let $V,H$ be separable Hilbert spaces with a compact imbedding into $V\hookrightarrow H$. Let $(\cdot,\cdot)$ be the inner product in $H$. \ac{$=L^2(\Omega)$??}
Let the dual of $V$ be $V'$, with the duality pairing $\langle \cdot, \cdot \rangle$ being an extension of $(\cdot,\cdot)$. 
$a_{\bsy}(\cdot,\cdot)$ and $b_\bsy(\cdot,\cdot)$  are Hermitian, bounded sesquilinear forms on $V\times V$ with $a_\bsy(\cdot,\cdot)$ being uniformly strongly coercive. Moreover, let $\calF : \calU \to V', \bsy \mapsto \calF_{\bsy} $ be a parametrization of dual forms.
%, corresponding to the prestress source problem in \eqref{eq:sourceproblem_weak_star}.
The following results can be applied to the situation in Problem \ref{problem:transformed_linearbuckling}.

Let $\mathcal{A}_{\bsy} ,\mathcal{M}_{\bsy} : V \to V'$ \ac{$V' \to V$ in both cases} be  operators such that 
\begin{align}
\label{eq:DefAy}
a_{\bsy}(\bsu, \bsv) &= \langle \mathcal{A}_{\bsy} \bsu ,  \bsv \rangle, \\
\label{eq:DefMy}
b_{\bsy}(\bsu, \bsv) &= \langle \mathcal{M}_{\bsy} \bsu ,  \bsv \rangle,
%\mathcal{F}_{\bsy}(\bsu) &= \int_{\Omega} \tilde F_i u_i \dxi  + \oint_{\partial \Omega_T} \tilde T_i u_i \mathop{dS},
\end{align}
for all $\bsu,\bsv \in V$ and $\bsy \in \calU$. Its well known that $\calA_\bsy$ and $\calA_\bsy^{-1}$ are homeomorphism from $V$ and $V'$. Compare, e.g. \cite[(3.2)]{DaugeandSuriNumericalapproximation}. %where $a_\bsy, b_{\bsy}$
%,\tilde F_i$ and $\tilde T_i$ are defined in \eqref{eq:bilinearforms},\eqref{eq:bilinearforms2}, \eqref{eq:transformedbodyload} and \eqref{eq:transformedtractionload}, respectively.
Moreover, we define 
\begin{align}\label{eq:definitionTy}
    \calT_{\bsy} = \calA_{\bsy}^{-1} \calM_{\bsy},
\end{align}
and then Problem \ref{problem:linearbuckling} is equivalent to solving for the largest eigenvalue in
\begin{equation}\label{eq:eigenvalueproblemTy}
    \calT_{\bsy} u(\bsy) = \mu(\bsy) u(\bsy), \quad \text{where} \quad \mu(\bsy) = \frac{1}{\lambda(\bsy)}.
\end{equation}
For the proofs of Lemma \ref{lemma:selfadjointedness} and \ref{lemma:pollutionfree}, see \cite[Section 3]{DaugeandSuriNumericalapproximation}: 
\begin{lemma}\label{lemma:selfadjointedness}
    $\calT_\bsy$ is a self-adjoint operator on $V$ with inner product $a_\bsy$.
\end{lemma}

We define the discrete and the essential spectrum
\begin{align*} 
    \sigma_{\ds}(\calT_\bsy) &:= \{ z \in \mathbb{C} \mid \ker( \calT_\bsy - z \mathbf{I}) \neq \{ 0 \},  \calT_\bsy - z \mathbf{I} \text{ is a Fredholm operator from $V$ into $V$ }  \}, \\
    \sigma_{\ess}(\calT_\bsy) &:= \{ z \in \mathbb{C} \mid \calT_\bsy -z \mathbf{I} \text{ is not a Fredholm operator from $V$ into $V$} \}.
    \end{align*}
As indicated in Remark \ref{remark:noncompactOperator}, $\calT_\bsy$ can potentially be a family of non-compact operators. However, we may assume that the essential spectrum is bounded in 
\begin{equation*}
    \sigma_{\ess}(\calT_\bsy) \subseteq [-D,D]
\end{equation*}
where $D > 0 $ is independent of $\bsy$. Note that, in the linear buckling model, $D$ depends only on a uniform Korn constant $K > 0$ and a uniform bound on the presstress norm $\| \bssigma_* (\bsy) \|_{L^\infty(\Omega; \mathbb{T}_2)} \leq M $, that is, $D = K M$. See \cite[Theorem 5.1]{DaugeandSuriNumericalapproximation}.
%This assumption is justified by the following result in \cite{DaugeandSuriNumericalapproximation} and a related result on the critical buckling load for thin domains.
Assuming the boundedness of the essential spectrum of $\calT_\bsy$ is justified by the results in the following remark:
\begin{lemma}[Pollution-free Eigenvalue Problem for Thin Domains]\label{lemma:pollutionfree}
    If $t > 0$ is the thickness of the domain, then the eigenvalues of Problem \ref{problem:linearbuckling} decay as $\calO(t^2)$, \cite{dauge:hal-00012083}.  Hence, $ \mu_1(\bsy)$ converges quickly towards infinity.
    Moreover, the interval $[-D,D]$ is independent of $t$, and thus the largest eigenvalues of $\calT_\bsy$ can be approximated without spectral pollution. That is, eigenvalues of the discrete solver do not seem to converge, since there are part of the essential spectrum of $\calT_\bsy$.
    %enabling a pollution-free eigenvalue
    Therefore, in the remainder, we assume that $t$ is small enough.
\end{lemma}


%The smallest buckling load $\lambda_1(\bsy)$ corresponds to the largest eigenvalue $\mu_1(\bsy)$ of $\calT_{\bsy}$. 
With the latter result, we may assume that the largest eigenvalues of $\calT_\bsy$ can be written in descending order
\begin{align*}
    \mu_1(\bsy) \geq \mu_2(\bsy) \geq  \mu_3(\bsy) \geq \ldots 
\end{align*}
and corresponding eigenvectors $u_1(\bsy), u_2(\bsy),\ldots $.
Moreover, we denote the eigenspace up to $k$ as 
\begin{equation}\label{eq:eigenspaceU}
    U_k(\bsy) = \mathrm{span}\{  u_i(\bsy) \mid i=1,\ldots,k \} .
\end{equation}




\subsection{Notation For Gevrey-Regularity}
The following definitions are from  \cite{chernov2026gevreyclassregularitysteadystate}. 

For $m \in \NN_0$ the falling factorial $\brac{\tfrac{1}{2}}_m$ is defined by  
\begin{equation*}
	\brac{\tfrac{1}{2}}_m
	:=
	\begin{cases}
		1 \quad & \text{if } m=0, \\
		\displaystyle \prod_{{k=0}}^{m-1} \brac{\tfrac{1}{2} - k}& \text{if } m > 0. 
	\end{cases}	
\end{equation*}	
 We introduce the quantity
\begin{equation*}
	\gdnota{m} := \abs{\brac{\tfrac{1}{2}}_m}.
\end{equation*}	
To streamline the presentation, we introduce the shorthand notation
	\begin{equation*}
		\fallfac{\delta}{\vnu} :=  \gdnota{\abs{\vnu}}  (\abs{\vnu}!)^{\delta-1}, {\qquad \vnu \in \cF.}
	\end{equation*}
The set of finitely supported sequences of non-negative integers $\cF$ is defined as
\begin{equation}\label{cF-def}
	\cF := 
	\left\{
	\vnu ={(\nu_j)_{j \in \NN}} : \nu_j\in \NN_0 \text{ with } \nu_j \neq 0 \text{ for  finitely many } j
	\right\}.
\end{equation}
For multi-indices $\valpha, \vbeta \in \cF$, the algebraic operation $\valpha + \vbeta$ and the partial orderings $\valpha \leq \vbeta$ are understood component-wise. In particular, $\valpha < \vbeta$ means $\valpha \leq \vbeta$ and $\valpha \neq \vbeta$. For a given sequence of positive real numbers $\vR=(\rho_j)_{j{\in \NN}}$ and a multi-index $\vnu \in \cF$, it is convenient to define
\begin{align*}
	\vR^{\vnu} := \prod_{j{\in \NN}} \rho_j^{\nu_j},
	\qquad \qquad
	\abs{\vnu} := \sum_{j{\in \NN}} \nu_j, 
	\qquad \qquad
	\vnu! := \prod_{j{\in \NN}} \nu_j!. 
\end{align*}


For a $\vnu\in \cF$ with  $\supp (\vnu) \subset\sett{1,2,\dots, n}$ and a function $v$ on $\R^\NN$, the partial derivative $\partial^{\vnu} v$ with respect to the variables $\vy$ are defined as
\begin{align*}
	\partial^{\vnu} v
	=
	\frac{\partial^{\abs{\vnu}}v}
	{\partial y_1^{\nu_1} \partial y_2^{\nu_2} \dots \partial y_n^{\nu_n}}
\end{align*}
with the convention that $\partial y_j^{\nu_j}$ is omitted if $\nu_j =0$. 


 The following definition of Gevrey-$\delta$ functions with countably many parameters will be used in our regularity analysis.
\begin{definition}[Gevrey-Regularity] \label{def:G-delta-def} 
	Let $\delta \geq 1$, $B$ be a Banach space, $I \subset \R^\NN$ be an open domain and a function $f:I \to B$ be such that its $\vy$-derivatives $\partial^{\vnu} f:I \to B$ are continuous for all $\vnu \in \cF$. We say that the function $f$ is of class Gevrey-$\delta$ if for each $y_0 \in I$ there exist an open set $J\subseteq I$ that contains $y_0$, and strictly positive constants $\vR = (\rho_1,\rho_2,\dots) \subset \R_{>0}^{\NN}$ and $C \in \R_{>0}$ that the derivatives of $f$ satisfy the bounds
	\begin{equation}\label{G-delta-def}
		\|\partial^{\vnu} f(\vy)\|_B \leq \frac{C}{\vR^{\vnu}} (|\vnu|!)^\delta, \qquad \forall \vy \in J, \quad \forall \vnu \in \cF.
	\end{equation}
	In this case we write $f \in G^\delta(I,B)$.
\end{definition}


\subsection{A uniform Spectral Gap}\label{subsec:UniformSpectralGap}
%In this section, we assume that the parametrization on $\calT_\bsy : V\to V$ is Gevrey-regular. This statement will be shown in Section \ref{subsec:proofTyRegularity} in the context of random domain mappings and the linear buckling problem. The aim of this section is to determine a condition under which the largest eigenvalue cluster of $\calT_\bsy$ is isolated. Consequently, we prove that the spectral projection onto this isolated eigenspace is Gevrey-regular.


%%% Rewritting the paragraph
Eigenvalues of smoothly param   etrized linear operators are known to admit non-smooth dependencies, that is, Gevrey-regularity of $\calT_\bsy$ does not imply Gevrey-regularity of $\mu_j(\bsy)$. %does not need to be Gevrye-regular% $C^\infty$.
%That is, there are counterexamples where we fail to bound $\vert \D^\bsnu \mu_j(\bsy) \vert$.
Therefore, analyzing the regularity of the solution of the eigenvalue problem requires us to redefine the problem; find the smallest isolated eigenvalue cluster that contains the largest eigenvalue.
Here, an isolated eigenvalue cluster is a subset $J \subseteq \N$ and the corresponding subset of eigenvalues  $S_J(\bsy) = \{ \mu_i(\bsy)\mid i \in J \}$ such that there exists a constant $C>0$ and 
 \begin{align*}
     \textrm{dist}( S_J(\bsy), \sigma(\calT_\bsy)\setminus S_J(\bsy))  \geq C
 \end{align*} 
  for all $\bsy \in \calU$. The smallest isolated eigenvalue cluster containing $\mu_1(\bsy)$ is characterized by the smallest $k \in \N$ such that $\{1,\ldots,k\}$ is isolated and for any $r  <k $ and $ J = \{1,\ldots, r\} $ for any $r  <k $, we have \begin{align*}
      \textrm{dist}( S_{J}(\bsy), \sigma(\calT_\bsy)\setminus S_{J}(\bsy))=0.
  \end{align*}
  
 We show that the corresponding spectral operators of $S_J(\bsy)$ are regular.
 Here, spectral operators of a cluster of eigenvalues $S_J(\bsy)$ are the operators $P(\bsy) $ that projects $V $ onto the eigenspace $U_J(\bsy) $ and $\Lambda(\bsy) = \calT_{\bsy} P(\bsy) $.  A formal definition in terms of the resolvent of $\calT_{\bsy}$ is given below in \eqref{eq:spectralProjector} and \eqref{eq:CauchyIntegralFormulaForD}.
 In the following section, we focus on establishing a condition for the largest eigenvalue cluster to be isolated.   For this result, we only require the assumption that $\calT_\bsy$ is Gevrey-regular stated in Assumption \ref{ass:gevreyTy}.
%Provided that the eigenvalue cluster of $\mu_1(\bsy)$ 
 %corresponding to $ \{ \mu_1(0), \ldots, \mu_s(0) \}$, where $s$ is the algebraic multiplicity of $\mu_i(0)$ 
%falling factorial 
 
%Eigenvalues of parametrized linear operators are known to admit non-smooth dependencies, regardless of the regularity of the parameter-operator dependency. 
%This is a simple consequence of the fact that algebraic eigenvalues can be characterized by their characteristic polynomial. 
%To fix ideas, we consider the characteristic polynomial of $\calT_\bsy$ restricted on a finite dimensional subspace. Let $P(\bsy)$ be an orthogonal projection onto a finite dimensional $\calT_\bsy$-invariant subspace of constant dimension. Then
%\begin{equation}\label{eq:characteristicpoly}
 %   \chi(z) =\det( P^*(\bsy)\calT_\bsy P(\bsy) - z\mathbf{Id}_n) = \sum_{i=0}^n a_i(\bsy)z^{i} = \prod_{i=1}^n (z -\mu_i(\bsy)).
%\end{equation} 
%The regularity of the coefficients $a_i(\bsy)$ does not imply that the roots are regular with respect to $\bsy$. 
%However, by comparing coefficients  $a_i$ and the products of $\mu_i(\bsy)$ in \eqref{eq:characteristicpoly}, we can exactly characterize which quantities remain regular. For example, the trace is \begin{align*}
%\mathrm{tr}(P^*(\bsy)\calT_\bsy P(\bsy)) = \sum_{i=1}^n \mu_i(\bsy) = -a_{n-1}(\bsy)
%\end{align*} 
%    and therefore remains regular.
%Moreover, the $\calT_\bsy$-invariant subspace is also an isolated subspace, because any crossing of non-zero invariant subspaces would increase its dimension.
%Hence, we need to find isolated invariant-subspaces or their corresponding eigenvalue clusters. 
%The criteria that we establish in the following,  uses the regularity constants of $\calT_\bsy$ that are shown in Section \ref{sec:regularityOfTy}. Numerical evidence in Section \ref{sec:numerical-example} suggests that the smallest eigenvalue cluster associated with the critical buckling load. 


We define the norm$\| \cdot \|_{\text{op}}$ as the operator norm for linear operators linear operators on $V \to V$. For the linear buckling model \eqref{eq:eigenvaluebucklingproblem}, the following assumption is shown in Section \ref{sec:regularityOfTy} under some relatively mild assumptions on the random domain mapping $\Phi$.
\begin{assumption}\label{ass:gevreyTy} 
For some $\delta \geq 1$, we assume that $\calT_\bsy$ is Gevrey-$\delta$-regular, that is,
\begin{equation}\label{eq:assumptionTyregular}
            \| \Dnuy \calT_\bsy \|_{\op} \leq C_{\calT} \bsb^{\bsnu} \fallennudelta, \quad  \text{ for all } \bsnu \in \calF, \bsy \in \calU.
 \end{equation}
Moreover, we assume that $\bsb \in \ell^1(\N)$
and $\mu_1(\bsy)  > D$ for all $\bsy \in \calU$. Hence, the largest eigenvalue is bounded away from the essential spectrum of $\calT_\bsy$.    
\end{assumption}



We can prove a Lipschitz property using the mean-value theorem to bound the variation of $\calT_{\bsy}$. 
\begin{lemma}[Lipschitz constant of $\calT_{\bsy}$]\label{lemma:variationboundTy}
    For any $ \bsy, \bsy_0 \in \calU$, we have
    \begin{align}\label{eq:boundvariationTy}
        \| \calT_{\bsy} - T_{\bsy_0} \|_{\op} \leq r  C_{\calT}\| \bsb \|_{\ell^1(\N)} =: C_V(r,C_{\calT},\bsb)= C_V,
    \end{align}
   \ac{$C$ is $\frac{1}{2}C_\calT$?} where $r = \text{diam}(\calU)$ and $C_{\calT},\rho_{\calT}$ are defined in Lemma \ref{lemma:GevreyRegularityTy}.
\end{lemma}
\begin{proof}
    Note that, 
   \begin{align*}
        \sum_{i=1}^\infty  \| \D^{e_i}_{\bsy} \calT_{\bsy} \|_{\op} \leq  C_{\calT}\sum_{i=1}^\infty  b_i =  C_{\calT} \| \bsb \|_{\ell^1(\N)} < \infty.
   \end{align*} 
   Since $\calU$ is convex, the mean-value theorem 
  implies that there is some $\zeta \in \calU$ between $\bsy_0$ and $\bsy$ such that 
   \begin{align*}
       \calT_{\bsy} - \calT_{\bsy_0} = \sum_{i=1}^\infty \D^{e_i}_{\bsy} \calT_{\zeta} (\zeta_i - \bsy_{0,i})
   \end{align*}
   and thus
   \begin{align*}
       \| \calT_{\bsy} - \calT_{\bsy_0} \|_{\op}  \leq  \sum_{i=1}^\infty  \| \D^{e_i}_{\bsy} \calT_{\zeta} \|_{\op}  \| \zeta - \bsy_0 \|_{\calU} 
        \leq  r C_{\calT}\| \bsb \|_{\ell^1(\N)}.
   \end{align*}
   \ac{$\calU$-norm is not defined, $\zeta$ should be boldface.}
\end{proof}
The perturbation of eigenvalues can be bounded by the same constant on the variation of the operator $\calT_\bsy$. This goes back to Weyl and a proof is included for convenience and because the special treatment of discrete and essential spectrum has not been found in the literature.
\begin{lemma}[Weyl's inequality]\label{lemma:Weylsinequality}
Let $A$ and $B$ be a self-adjoint operators such that for $M \in \{A,B\} $, $\sigma_{\ds} (M)\cap (\R \setminus \sigma_{\ess}(M))$ is not empty and 
let $\mu_1(M) \geq \mu_2(M)  \geq \ldots $ denote the eigenvalues  for which $\mu_k(M) \in \sigma_{\mathrm{ds}} (M)\cap (\R \setminus \sigma_{\ess}(M))$ in descending order. 
Then
\begin{align*}
    \mu_{i +j -1}(A+ B) \leq \mu_i(A) +  \mu_{j}(B),
\end{align*}
for all $i,j \geq 1$.
Consequently,
\begin{align*}
    \vert \mu_{k} (\bsy)- \mu_k(\bsy_0) \vert \leq C_V, \quad  \text{ for all } \bsy \in \calU,
\end{align*}
where $V$ is defined in \eqref{eq:boundvariationTy}, for all $k  \geq 1$ such that the eigenvalues $\mu_k(\bsy) ,\mu_k(\bsy_0) $ are not in numerical range of the essential spectrum of $\calT_\bsy$.
\end{lemma}
\begin{proof}
    Weyl's inequality is 
    \begin{equation*}
        \mu_{i+j-1}(A+ B)  \leq \mu_{i}(A) + \mu_j(B), \quad i,j \geq 1.
    \end{equation*}
   % for operators $A,B$ that satisfies the with a non-empty discrete spectrum, where the essential spectrum is bounded to some interval around $0$. At least for the discrete spectrum, that is not mixed with the essential spectrum, one can prove above inequality.
    This can be shown by using the $\min$-$\max$ characterization for self-adjoint operators and the identity holds for the discrete \ac{what do you mean by algebraic?} spectrum
    
    \begin{align*}
        \mu_{i} (A) = \min_{ \dim(W) \leq i-1} \max_{\substack{u \bot W \\\| u \| = 1}}  \langle Au,u \rangle.
    \end{align*}
    Let $S_A$ and $S_B$ be the subspaces of dimension $i-1$ and $j-1$, respectively, for which the minimum is attained.
    Then $S = S_A + S_B$ has dimension less than $i+j-1$. Furthermore, $u \bot S$ implies $u \bot S_1$ and $u \bot S_2$.
    Thus
 \begin{align*}
     \mu_{i+j-1} (A + B) &= \min_{ \dim(W) \leq i+j-2} \max_{\substack{u \bot W \\\| u \| = 1}} \, \langle (A+B)u,u \rangle \\
     &\leq \max_{ \substack{u \bot S \\\| u \| = 1}} \, \langle Au,u \rangle + \langle Bu,u \rangle \\
     &\leq \mu_{i}(A) + \mu_j(B).
 \end{align*}
 For the second statement, we set $j = 1$, and $A = \calT_{\bsy}, B = \calT_{\bsy_0}-  \calT_{\bsy}$ and immediately obtain 
\begin{align*}
     \mu_{k} (\bsy)- \mu_k(\bsy_0)  \leq \mu_1(\calT_{\bsy_0}-\calT_{\bsy})  \leq \| \calT_{\bsy_0}-\calT_{\bsy} \|_{\op}  \leq C_V.
\end{align*} Interchanging $\bsy_0$ and $\bsy$, yields the claim.
\end{proof}


\begin{lemma}[Uniform Spectral Gap]\label{lemma:uniformSpectralGap}
    Let $C_V$ the bound on the variation of $\calT_\bsy$ from \eqref{eq:boundvariationTy}. %Moreover,  assume that $s \in \N$ is the algebraic multiplicity of $\mu_1(0)$. Further, assume that the essential spectrum is sufficiently far away from $\mu_1(0)$, that is,  
    %\begin{align*}
    %\mu_1(0) - V > D,
    %\end{align*} where $D$ is such that $\sigma_{\ess}(\calT_\bsy) \subseteq [-D,D]$.
    Let  $s \in \N$ be minimal such that 
    \begin{equation}\label{eq:def_delta0}
        \delta_0 := \mu_{s}(0)-\mu_{s+1}(0)  > 0. 
    \end{equation} 
    If 
    \begin{equation}\label{eq:def_delta_hat}
        \hat \delta = \delta_0 -2C_V > 0,
    \end{equation} then the eigenspace of the first $s$ eigenvalues $U_s(\bsy)$, see \eqref{eq:eigenspaceU}, is isolated from the rest of the spectrum with uniform spectral gap that is at least of size $\hat \delta$.
\end{lemma}
\begin{proof}
%First, note that for $k \in \{1,\ldots,s\}$, 
%\begin{align*}
%    \mu_k(\bsy) - D =  \mu_k(\bsy)  -  \mu_k(0) +  \mu_k(0) -D  \geq   \mu_1(0) - V - D > 0
%\end{align*}
%which implies that $\mu_k(\bsy) \notin [-D,D]$ for all $\bsy \in \calU$. Hence there is no mixing of the essential spectrum with the first eigenvalue cluster.
Let $i \geq 1$ and $\bsy \in \calU$. 
    Then 
    \begin{align*}
        \mu_s(\bsy) - \mu_{s+i}(\bsy)  & = \underbrace{\mu_s(\bsy)- \mu_s(0)}_{\geq -C_V } + \underbrace{\mu_s(0) - \mu_{s+i}(0)}_{\geq \delta_0}  + \underbrace{\mu_{s+i}(0) - \mu_{s+i}(\bsy)}_{\geq -C_V} \\
        &\geq \delta_0 -2C_V = \hat \delta.
    \end{align*}
\end{proof}
\subsection{Gevrey-Regularity of an Isolated Eigenvalue Cluster}\label{subsec:clusterRegular}
Assuming that the conditions in Lemma \ref{lemma:uniformSpectralGap} are satisfied. We can define the resolvent and a contour line $\Gamma$ that separates $\{ \mu_1(\bsy), \ldots, \mu_s(\bsy) \}$ from the rest of the spectrum such that 
\begin{equation*}
    z \mapsto R(z,\bsy):= (\calT_\bsy-z\mathbf{I}))^{-1}
\end{equation*} 
is analytic for all $z \in \Gamma$ and $\bsy \in \calU$. 
%The aim of this section is to prove Theorem \ref{theorem:main}.
%The proof avoids the direct differentiation of the buckling modes which is ill-posed, because potential crossings of eigenmodes makes their $\bsy$-dependence non-smooth.
%And instead the proof relies on the resolvent formalism and Cauchy integrals.
%Based on Theorem \ref{theorem:main1}, we have an explicit criterion on when we can expect to find an eigenspace isolating curve $\Gamma$.
In such a case, the Cauchy integral formula for the resolvent can be applied \cite[Chapter III, \S4 , Theorem 6.17]{kato1995perturbation}. 
%exhibiting the enclosed eigenspace in form of the spectral projection operator $P(\bsy)$ in \eqref{eq:spectralProjector}, and the resolvent operator, restricted on the enclosed eigenvalues $\mu_1(\bsy), \ldots, \mu_s(\bsy)$. 
%Here, $s$ is the size of the first eigenvalue cluster, as in \eqref{eq:choiceofs}. 
That is, the operators
\begin{align}\label{eq:spectralProjector}
    P(\bsy) &:= \frac 1 {2\pi i} \int_{\Gamma} (\calT_{\bsy} -z \mathbf{I})^{-1}  \mathop{dz}, \\
    \label{eq:CauchyIntegralFormulaForD}
    \Lambda(\bsy) &:=  \frac  1 {2\pi i} \int_{\Gamma} z \mathbf{I} (\calT_{\bsy} - z \mathbf{I})^{-1}\mathop{dz}  
\end{align}
   % are Gevrey-$\delta$-regular on $\calU$.
are well-defined and moreover, 
\begin{equation*}
    \mathrm{ran} (P(\bsy)) = U_s(\bsy), \quad P^2(\bsy) = P(\bsy),
\end{equation*} where $U_s(\bsy)$ is defined in \eqref{eq:eigenspaceU} and 
\begin{equation*}
    \Lambda(\bsy) =\calT_{\bsy}P(\bsy),
\end{equation*}
for all $\bsy \in \calU$.

Then the trace operator of $\Lambda(\bsy)$ yields
\begin{equation}
    \mathrm{tr} \,\Lambda(\bsy) =  \mu_1(\bsy)+  \ldots + \mu_s(\bsy),
\end{equation}

%while the projector $P_\bsy = P(\bsy)$ in \eqref{eq:spectralProjector} satisfies $P_\bsy (\calL(V)) = U_s(\bsy)$. Here, recall that $ U_s(\bsy)$ is the eigenspace, spanned by the first $s$ eigenvectors of $\calT_\bsy$. In other words, we have

%The resolvent theory is, for example, explained in great detail in Kato \cite[Chapter I, \S 3 \& Chapter III, \S 4]{kato1995perturbation}. Note that, it is sufficient that $\calT_{\bsy}$ is linear and bounded.
%Define the resolvent operator of $\calT_{\bsy}$ %for $z \in \mathbb{C}$ and $\bsy \in \calU$:
%\begin{equation}
    %R(z,\bsy) = ( \calT_{\bsy} - z )^{-1}.
%\end{equation}
%All linear and bounded operators are closed operators. 
%Thus by the analysis in \cite[Chapter I, \S 3]{kato1995perturbation}, the resolvent $z \mapsto R(z,\bsy)$ is a meromorphic function, and the Cauchy-integrals $P(\bsy), \Lambda(\bsy)$  in \eqref{eq:spectralProjector} and \eqref{eq:CauchyIntegralFormulaForD} are well-defined for all $\bsy \in \calU$.  

%Let $\Gamma \subset \mathbb{C}$ be a fixed closed contour isolating $\mu_1(\bsy)$ from the rest of the spectrum $\sigma(\calT_{\bsy})$. Due to the uniform spectral gap, such a contour exists and the resolvent $R(z,\bsy)$ is uniformly bounded for all $z \in \Gamma$ and $\bsy \in \calU$.
%\begin{theorem}
%Under Assumptions \ref{subsection:assumptions}, the critical buckling load $\lambda_1(\bsy)\in G^\delta(\calU,\RR)$, i.e. there exist constants $C > 0$ and a radius sequence $\vR$ such that for any $\vnu\in\cF$:
%\begin{equation}
    %|\partial^\nu_y \lambda_1(\bsy)| \le C \frac{(|\vnu|!)^\delta}{\vR^{\vnu}}.
%\end{equation}
%\end{theorem}

\begin{theorem}(Gevrey-Regularity of an isolated Eigenvalue Cluster)\label{theorem:ClusterRegularity}
    Assume that Assumption \ref{ass:gevreyTy} and the conditions in  in Lemma \ref{lemma:uniformSpectralGap} are satisfied. 
    %That is, $\calT_{\bsy}$ is Gevrey-regular and there is a uniform spectral gap that separates $\{ \mu_1(\bsy),\ldots, \mu_s(\bsy) \}$ from the rest of the spectrum for all $\bsy \in \calU$. 
    The contour
    \begin{equation}\label{eq:contourDef}
    \Gamma = \partial B( \mu_1(0), \delta_0/2) \subseteq \mathbb{C},
\end{equation} where $\delta_0 $ is the initial gap defined in \eqref{eq:def_delta0}, maximizes the distance between $\Gamma$ and $\sigma(\calT_\bsy) $ for all $\bsy \in \calU$. This choice minimizes the bound on the resolvent, that is, 
\begin{align}\label{eq:resolventUniformbound}
    \| R(z,\bsy) \|_{\op} \leq \gamma^{-1}, \quad \text{ for all } \bsy \in \calU, z \in \Gamma.
\end{align}
where $\gamma = \hat \delta /2 $ and $\hat \delta $ is defined in \eqref{eq:def_delta_hat}. Then $P(\bsy)$ and $\Lambda(\bsy)$ 
    %defined in \eqref{eq:spectralProjector} and \eqref{eq:CauchyIntegralFormulaForD} 
    are Gevrey-regular with
    \begin{align*}
		\norm{\partial_{\vy}^{\vnu} P(\vy)}{\op}
		&\leq \delta_0 C_R \rho_R^{\vert \bsnu \vert} \bsb^{\bsnu} \fallennudelta, \\
  \norm{\partial_{\vy}^{\vnu} \Lambda(\vy)}{\op}
		&\leq \widetilde C_\Gamma  \delta_0 C_R \rho_R^{\abs{\vnu}} \vb^\vnu \fallennudelta.
	\end{align*}
    for all $\bsnu \in \calF$ and $\bsy \in \calU$, where the constants are defined as
    \begin{align} 
       \widetilde C_\Gamma &= \mu_1(0) + \delta_0/2, \\
         C_{\calT_\bsy,\Gamma}&=   \max \{ \gamma^{-1}, C_\calT \}, \label{eq:CcalTGamma} \\
   C_R &= C_{\calT_\bsy,\Gamma}/3,  \label{eq:C_R}\\
   \rho_R &= 3  C_{\calT_\bsy,\Gamma}^2   \label{eq:rho_R} , 
    \end{align}
    where $\delta_0$ is defined in \eqref{eq:def_delta0} and $C_\calT $ in \eqref{eq:assumptionTyregular}.
\end{theorem}
To the best of our knowledge, there is no previous work on this or similiar statements, however, a comprehensive overview on related resolvent techniques are covered in the textbook by Kato \cite{kato1995perturbation}. Key concepts are required, such as, Lemma \ref{lemma:productInversionrule} which is based on novel developments in \cite{chernov2026gevreyclassregularitysteadystate}.
\begin{proof}
%The proof avoids the direct differentiation of the buckling modes, which is ill-posed in a 2D eigenspace, and instead relies on the resolvent formalism and Cauchy integrals.

%\paragraph{Step 1: Operator Formulation.}
%We rewrite \eqref{eq:variational_buckling} as a standard operator eigenvalue problem. Let $\calT_{\bsy} : V \to V$ be the compact operator defined by $\calT_{\bsy} = \calA_{\bsy}^{-1} \calM_{\bsy}$. The problem becomes
%\begin{equation}
%    \calT_{\bsy} D(\bsy) = \mD(\bsy) D(\bsy), \quad \text{where} \quad \mD(\bsy) = \frac{1}{\lambda(\bsy)}.
%\end{equation}
%The smallest buckling load $\lambda_1(\bsy)$ corresponds to the largest eigenvalue $\mu_1(\bsy)$ of $\calT_{\bsy}$. By Assumption \ref{assump: non-splitting}, $\mu_1(\bsy)$ has multiplicity 2 for all $\bsy \in \calU$.

%\paragraph{Step 2: The Resolvent and Cauchy Integral.}
%Define the resolvent operator of $\calT_{\bsy}$ for $z \in \mathbb{C}$:
%\begin{equation}
 %   R(z,\bsy) = (zI - \calT_{\bsy})^{-1}.
%\end{equation}

 % \hat \delta /2 = (\delta_0-V)/2$. 
%Furthermore, let $\Gamma \subset \mathbb{C}$ be the contour isolating $\mu_1(\bsy), \ldots, \mu_s(\bsy)$ from the rest of the spectrum $\Sigma(\calT_{\bsy})$ for all $ \bsy \in \calU$.
%
We start by showing that the resolvent $R(z,\bsy)$ is uniformly bounded as in \eqref{eq:resolventUniformbound} with the choice of $\Gamma$ defined in  \eqref{eq:contourDef}. This property is straight forward consequence from the fact that $\mathrm{dist}(\Gamma, \sigma(\calT_\bsy)) \geq \gamma$ for all $\bsy \in \calU$.

To see this, $\Gamma$ intersects the real line in the center of $\mu_1(0)$ and $\mu_{s+1}(0)$ at $m := \mu_1(0) - \delta_0/2  = \mu_{s+1}(0) +  \delta_0/2 $ and since the variation of both eigenvalues are bounded by $C_V$ and other eigenvalues are smaller than $\mu_{s+1}(0) + C_V$, the distance between all eigenvalues and $m$ is at least $(\delta_0 -2C_V)/2 = \gamma$. 

We show this step more formally:  \tododuc{is this really needed?}
%More precisely, we show $\mathrm{dist}(\Gamma, \sigma(\calT_\bsy) \geq \gamma$, where $\hat \delta $ is defined in \eqref{eq:def_delta_hat}. 
Let $\bsy \in \calU$ and $i \in \N^*$, such that $\mu_{i}(\bsy) \in \sigma_{\mathrm{ds}}(\calT_\bsy)$ and $ z \in \Gamma$. 
Then first assume that $\mu_i(0) \neq \mu_1(0)$, that is, $\mu_i(\bsy)$ is not in the first eigenvalue cluster $ \{\mu_1(\bsy),\ldots, \mu_s(\bsy) \}$, and thus
\begin{align*}
         \vert \mu_{i}(\bsy) - z  \vert &=  \vert  \mu_{i}(\bsy) -  \mu_{i}(0) + \mu_{i}(0) -  \mu_{1}(0) +\mu_{1}(0)  - z\vert  \\ 
     %    &\geq \vert   \mu_{i}(0) -  \mu_{1}(0) \vert - \vert \mu_{i}(\bsy) -  \mu_{i}(0)  + \mu_{1}(0)  - z \vert   \\
          &\geq \vert   \mu_{i}(0) -  \mu_{1}(0) \vert - (\vert \mu_{i}(\bsy) -  \mu_{i}(0) \vert   + \vert \mu_{1}(0)  - z \vert)   \\
         &\geq \delta_0  - C_V  - \delta_0/2= \gamma
\end{align*} where $C_V$ bounds the variation of $\mu_i(\bsy)$ as in Lemma \ref{lemma:Weylsinequality}. \ac{2nd last inequality unclear}
If  $\mu_i(0) = \mu_1(0)$, we have 
\begin{align*}
    \vert \mu_{i}(\bsy) - z  \vert   &= \vert  \mu_{i}(\bsy) -  \mu_{i}(0) +\mu_{1}(0)  - z\vert 
    \\ & \geq  \vert \mu_{1}(0)  -z \vert  -  \vert   \mu_{i}(\bsy) -  \mu_{i}(0)\vert \\
    &\geq \delta_0/2 -C_V = \gamma.
\end{align*}\ac{2nd last inequality unclear}
%We denote the spectral radius of a bounded linear operator $T $ as $\mathrm{spr}(T) = \liminf_{n\to \infty }  \| T^n \| ^{1/n}$. 
For a linear operator $T$, the spectral radius can be defined as \cite[I-(5.12)]{kato1995perturbation}
\begin{align*}
    \mathrm{spr}(T) = \max \{ \vert \lambda \vert : \lambda \in \sigma(T) \}.
\end{align*}
Since $\calT_\bsy$ is a self-adjoint operator on a Hilbert space $(V, \langle \cdot, \cdot\rangle)$, the basic result (e.g.  \cite[V-(2.4)]{kato1995perturbation})  $\mathrm{spr}  (\calT_\bsy) = \| \calT_\bsy \|_\op $ holds. Moreover, the spectrum of the residual is 
\begin{align*}
\sigma((T-\zeta)^{-1}) = 1/(\sigma(T)-\zeta), \quad \zeta \in \mathbb{C},
\end{align*}
which can be seen from $(T-\zeta)^{-1} - \lambda^{-1} =  \lambda^{-1} (T-\zeta)^{-1} (\lambda - (T-\zeta))$. 
And hence
\begin{align*}
     \| R(z,\bsy)  \|_{\op} =  \max\{1/\vert \lambda - z \vert : \lambda \in \sigma(\calT_\bsy) \}  =  \frac{1} {\text{dist}(\sigma(\calT_{\bsy}), z)}  \leq \gamma^{-1},
\end{align*}
for all $ \bsy \in \calU, z \in \Gamma$.

%where in the last equality the second term $(\mu - z) \langle u, u\rangle$ is independent of $u$ and the first term is equal to zero, because $\mu$ is in the spectrum. 
%also requires self-adjointness such that  $ \langle (\calT_{\bsy}-\mu \mathbf{I})u,u \rangle  \geq 0$ holds.\
\ac{2nd last inequality unclear}
%Consequently, the resolvent $R(z,\bsy)$ is uniformly bounded for all $z \in \Gamma$ and $\bsy \in \calU$.
%\begin{equation*}
%    \| R(z,\bsy) \|_{\op} \leq \gamma^{-1} \quad \text{ for all }  \bsy \in \calU, z \in \Gamma,
%\end{equation*} 
By Assumption \ref{ass:gevreyTy}, we have
\begin{equation*}
    \| \Dnuy  (\calT_\bsy - z )\|_{\op} =   \| \Dnuy  \calT_\bsy \|_{\op}\leq C \bsb^{\bsnu} \fallennudelta, \quad \forall \bsnu \in \calF, \bsy \in \calU.
\end{equation*}
Thus, using the inversion rule for linear and bounded operators in Lemma \ref{lemma:productInversionrule} with the constants $C_R$ and $\rho_R$ as in \eqref{eq:C_R} and \eqref{eq:rho_R}, we have 
\begin{equation*}
    \| \Dnuy R(z,\bsy) \|_{\op} \leq C_R \rho_R^{\vert \bsnu \vert} \bsb^{\bsnu} \fallennudelta
\end{equation*}
for all $\bsnu \in \calF$ and $\bsy \in \calU$ and $z \in \Gamma$.

Now, we are in the position to show that the integral operators $\Lambda(\bsy)$ and $P(\bsy)$ are Gevrey-regular. For this step, let $\bsy \in \calU$ be fixed.
Since $z \mapsto R(z,\bsy)$ and $ z \mapsto z \mathbf{I}  R(z,\bsy)$ are continuous over the $\bsy$-independent domain $\Gamma$\ac{$R$ should be also sufficiently regular in $z$, I suppose}, we can interchange differentiation and integration:
\begin{equation}
    \partial^\nu_y \Lambda(\bsy) = \frac{1}{2\pi i} \oint_{\Gamma} z \left( \partial^\nu R(z,\bsy) \right) dz,
\end{equation}
and
	\begin{equation*}
		\partial_{\vy}^{\vnu} P(\vy) = \frac{1}{2\pi i} \int_{\Gamma} \partial_{\vy}^{\vnu} R(z, \bsy)  \, dz.
	\end{equation*}
	%Taking the Frobenius norm on both sides and utilizing the submultiplicative property $\norm{A B}{F} \leq \norm{A}{F}\norm{B}{F}$ (as noted in \cite[ (14)]{chernov2026gevreyclassregularitysteadystate}), along with the dimensional factor $m \ell$, we obtain:
    We obtain
	\begin{align*}
		\norm{\partial_{\vy}^{\vnu} P(\vy)}{\op}
		&\leq \frac{\abs{\Gamma}}{2\pi} \sup_{z \in \Gamma} \norm{\partial_{\vy}^{\vnu} R(z, \vy)}{\op} \\
		&\leq \frac{\abs{\Gamma}}{2\pi} C_R \rho_R^{\vert \bsnu \vert} \bsb^{\bsnu} \fallennudelta .
	\end{align*}
	
	Similarly, for the integral $\Lambda(\vy)$, the presence of the factor $z$ inside the integrand introduces an additional multiplier bounded by $\widetilde C_\Gamma := \mu_1(0) + \delta_0/2  =  \max_{z \in \Gamma} \abs{z} $. Following the exact same procedure, we arrive at:
	\begin{align*}
		\norm{\partial_{\vy}^{\vnu} \Lambda(\vy)}{\op}
		&\leq \widetilde C_\Gamma  \frac{\abs{\Gamma}}{2\pi}C_R \rho_R^{\abs{\vnu}} \vb^\vnu \fallennudelta.
	\end{align*}
	With  $  \abs{\Gamma}/(2\pi) = \delta_0 $, this completes the proof.
\end{proof}

\subsection{Gevrey-regularity of the Riesz representers}
\label{subsec:proofTyRegularity}
We complete this section with novel results about the Gevrey-regularity of composition and inversion of bounded linear operators and the linear representers of bilinear forms. 
Combined with the standard result on the Gevrey-regularity of the coefficients $\tilde \calC^{ijkl}(\bsy)$ in \eqref{eq:transformed_stiffness} and $
\tilde T^i)(\bsy)$ in \eqref{eq:transformedtractionload}, we can show that $\calT_{\bsy}$ is indeed Gevrey-regular under Assumption~\ref{assumption:GevreyRegularityJ}.
%$\Phi$. Combined with this standard result, we can show the Gevrey-regularity of $\calT_{\bsy}$, where now $a_\bsy$ and $b_{\bsy}$ are defined by the linear elasticity random domain formulation in \eqref{eq:bilinearforms} and \eqref{eq:bilinearforms2}. 

Gevrey-regularity in the sense of the falling factorial $\fallennudelta$ is preserved under
multiplication and inversion of linear operators \cite{CL2024}. 
%This development led to new proof techniques for the classical regularity analysis. Notably, in \cite{HarbrechtSchwabImplicit}, using the implicit function theorem and  \cite{chernov2026gevreyclassregularitysteadystate}, using the falling factorial technique, the authors introduce novel techniques to show regularity of quantities of interest related to the solution of parametrized partial differential equations. 
\cite[Lemma~2]{chernov2026gevreyclassregularitysteadystate} states this property for parametrized matrices, but the arguments are the same for the infinite-dimensional case. As a result, Gevrey-regularity extends to linear and bilinear operators shown in the following two Lemmas.
%, we propose a novel approach to show Gevrey-regularity of composition and inverse bounded linear operators between Banach spaces, leveraging the falling factorial 
%technique further.
%, enabling to show Gevrey-regularity of compositions and inversions of Gevrey-regular linear operators. 
In the following Lemma, the result in (iii) is known by \cite[Lemma~2]{chernov2026gevreyclassregularitysteadystate} but restated here for convenience.
\begin{lemma}[Composition, inversion and determinant]\label{lemma:productInversionrule}
Let $\calU \subseteq \R^\infty$ be compact and bounded, and $V,W,Z$ be Banach spaces. Moreover, let $\calL(X,Y)$ denote the set of bounded linear operators from $X$ to $Y$.
Let $K\colon\calU\to\calL(W,V)$ and $M\colon\calU\to\calL(V,Z)$ be Gevrey-regular, i.e.\ there are
$C_K\ge0, C_M\ \geq  1$, $\bsb\in\R_{\ge0}^\infty$ and $\delta\ge1$ with
\begin{align*}
  \|\Dnuy K_\bfy\|_{\calL(W,V)} &\le C_K\,\bsb^{\bsnu}\,\fallennudelta,\\
  \|\Dnuy M_\bfy\|_{\calL(V,Z)} &\le C_M\,\bsb^{\bsnu}\,\fallennudelta, \\
  \|M_\bfy^{-1}\|_{\calL(Z,V)}&\le C_M
\end{align*}
for all $\bfy\in\calU$ and $\bsnu\in\calF$.
Then for all $\bfy\in\calU$ and $\bsnu\in\calF$:
\begin{enumerate}
  \item[(i)] \emph{Composition:}
    $\|\Dnuy(M \circ K)_\bfy\|_{\calL(W,Z)}\le 4\,C_MC_K\,\bsb^{\bsnu}\,\fallennudelta$.
  \item[(ii)] \emph{Inversion} (for $\bsnu > 0$):
    $\|\Dnuy M_\bfy^{-1}\|_{\calL(Z,V)}\le\frac{C_M}{3}\,(3C_M^2)^{|\bsnu|}\,\bsb^{\bsnu}\,\fallennudelta$.
  \item[(iii)] \emph{Determinant:} if $A\colon\calU\to\R^{d\times d}$ satisfies
    $\|\Dnuy A_\bfy\|\le C_A\,\bsb^{\bsnu}\,\fallennudelta$, then
    $|\Dnuy\det A_\bfy|\le d!\,(4C_A)^d\,\bsb^{\bsnu}\,\fallennudelta$.
    % CHECK: original had M^3_2 as target norm and no D^nu on A_y in the hypothesis.
\end{enumerate}
\end{lemma}
\textcolor{orange}{This is based on the general multivariate Leibniz rule and that part of the proof can be left out.}
\begin{proof}
\emph{(i) and (ii).}
We only need to show that the Leibniz rule holds, that is,
\begin{equation}\label{eq:LeibnizOperators}
  \Dnuy(M_\bfy \circ K_\bfy)=\sum_{\bseta\le\bsnu}\binom{\bsnu}{\bseta}\,
  \partial^{\bsnu-\bseta}_\bfy M_\bfy \circ \;\partial^{\bseta}_\bfy K_\bfy ,\qquad \bsnu\in\calF .
\end{equation}
Then taking the operator norms in \eqref{eq:LeibnizOperators} and argue exactly as in
\cite[Lemma~2]{chernov2026gevreyclassregularitysteadystate}, with operator norms in place of matrix norms yields both (i) and (ii).
We show the Leibniz rule by induction over $\vert \bsnu \vert$. 

\emph{Induction start}. We show the case $\vert \bsnu \vert = 1$. Fix one coordinate $t=y_j$. Adding and subtracting $M_{t+h}K_t$ and using the linearity of $M_{t+h}$ gives
\begin{align*}
  \frac{M_{t+h}K_{t+h}-M_tK_t}{h}
  &=M_{t+h}\,\frac{K_{t+h}-K_t}{h}+\frac{M_{t+h}-M_t}{h}\,K_t .
\end{align*}
Since $M$ and $K$ are differentiable in operator norm, $M_{t+h}\to M_t$ and the difference
quotients converge in operator norm. Hence
\begin{equation*}
  \partial_t(M_tK_t)=M_t\,\partial_tK_t+(\partial_tM_t)\,K_t .
\end{equation*}
\emph{Induction step}. Let $ \bsnu - w =e $ and  $e \geq 0 , \vert e \vert = 1$.
Then with Pascal's rule from above and the induction hypothesis on $w$, we have  
\begin{align*}
    \Dnuy (M_\bfy K_\bfy) &= \D^e_{\bfy} \left( \sum_{0 \leq \bseta \leq w} \binom{w}{\bseta} \D^{w-\bseta}_\bfy M_\bfy  \D^\bseta_\bfy K_\bfy\right) \\
    &=  \sum_{0 \leq \bseta \leq w} \binom{w}{\bseta} \left( \D^{w-\bseta +e}_\bfy M_\bfy  \D^\bseta_\bfy K_\bfy + \D^{w-\bseta}_\bfy M_\bfy  \D^{\bseta+e}_\bfy K_\bfy \right).
\end{align*}
Noting that $w  + e = \bsnu $ in the first term and replacing $\bseta = \bseta +e$ in the second term gives
\begin{align*}
    &= \sum_{0 \leq \bseta \leq w}    \binom{w}{\bseta} \D^{\bsnu -\bseta}_\bfy M_\bfy  \D^{\bseta}_\bfy K_\bfy  + \sum_{e \leq \bseta \leq \bsnu } \binom{w}{\bseta -e} \D^{\bsnu - \eta}_\bfy M_\bfy  \D^{\eta}_\bfy K_\bfy  \\
    &= \D^{\bsnu}_\bfy M_\bfy K_\bfy + \sum_{e \leq \bseta \leq w} \left(\binom{w}{\bseta} + \binom{w}{\bseta - e} \right) \D^{\bsnu - \bseta}_\bfy M_\bfy \D^{\bseta}_\bfy K_\bfy + M_\bfy \D^{\bsnu}_\bfy K_\bfy 
\end{align*}
Since only one component of  $\bseta - e $ and $ \bseta$ differ from each other, the standard formula $\binom{n-1}{k} + \binom{n-1}{k-1} = \binom{n}k$ can be applied to give $
\binom{w}{\bseta} + \binom{w}{\bseta - e}  = \binom{\bsnu}{\bseta}$. This concludes the proof.

\emph{(iii)} is proved as in \cite[Lemma~2]{chernov2026gevreyclassregularitysteadystate}.

\end{proof}
\begin{lemma} \label{lemma:RepresenterBilinearFormRegularity}
Let $\calU \subseteq \R^\infty$ be compact and bounded. Let $V, W$ be Banach spaces and let $a_\bsy: V \times W \to \R$ be a Gevrey-regular with $\delta \geq 1 $ such that 
 \begin{align}\label{eq:a_yGevrey}
      \| \D^\bsnu a_\bsy( \bsu, \bsv ) \|_{\calL(V \times W) } \leq C \bsb^{\bsnu} \fallennudelta =: C_{\bsnu}, \quad  \forall \bsnu \in \calF, \bsy \in \calU.
 \end{align}
Defining  $\calA_\bsy: V \to W'$ such that $a_\bsy( \bsu, \bsv) =  \langle \calA_\bsy \bsu, \bsv \rangle $ for $\bsu \in V$ and $\bsv \in W$, we the obtain that $\D^\bsnu \calA_\bsy$ exists in the weak sense for all $\bsnu \in \calF$ and $\calA_\bsy$ is Gevrey-regular with the same constants
  \begin{align*}
      \| \D^\bsnu \calA_\bsy \|_{\calL(V,W')} \leq C \bsb^{\bsnu} \fallennudelta, \quad  \forall \bsnu \in \calF, \bsy \in \calU.
  \end{align*}
\end{lemma}
\textcolor{orange}{I will try to find a reference for this. The result seems rather obvious.}
\begin{proof} Define $\calA_{\bsy}^{(\bsnu)}$ in the weak sense, via 
$  \langle \calA_{\bsy}^{(\bsnu)} u , v \rangle = \D^{\bsnu} a_{\bsy}(u,v)$  for $u \in V, v \in W$. Then 
\begin{align*}
    \langle  \calA_{\bsy}^{(\bsnu)} u, v \rangle  =  \D^{\bsnu} a_{\bsy}(u,v) \leq C_{\bsnu} \|u \|_V \|v\|_W,\quad \forall u \in V, v \in W,
\end{align*} by \eqref{eq:a_yGevrey} and  recursively, we have $\D^{\bsnu} \calA_{\bsy}u \to  \calA_{\bsy}^{(\bsnu)}  u$ in weak convergence for all $ u \in V$ and $ \bsnu \in \calF$.
    Then% we show that $\| \D^\bsnu \calA_\bsy \|_{\calL(V,W')} = \| \D^{\bsnu} a_\bsy \|_{\calL(V \times W) } $ if \eqref{eq:a_yGevrey} holds. 
    \begin{align*}
        \| \D^\bsnu \calA_\bsy \|_{\calL(V,W')} &= \sup _{ \| u \|_V= 1} \| \D^\bsnu \calA_\bsy u  \|_{\calL(W')} =  \sup _{ \| v \|_W = \| u \|_V= 1}  \langle \D^\bsnu \calA_\bsy u , v \rangle \\
        &=  \sup _{ \| v \|_W = \| u \|_V= 1}  \D^{\bsnu} a_{\bsy}(u,v)\\
        &= \| \D^{\bsnu} a_{\bsy} \|_{\calL(V \times W) } 
    \end{align*}
    where the second last equality follows from the weak convergence of $ \D^{\bsnu} A_{\bsy} u$.
 \end{proof}


\section{Regularity of Critical Buckling Load under Domain Uncertainty} \label{sec:regularityOfTy}
%In this section, we show that the random domain mapping approach leads to a Gevrey-regular spectral operator $\calT_\bsy$.


Under the assumption that the random domain mapping is Gevrey-regular and uniformly diffeomorphic, we show Gevrey-regularity of the operator $\calT_\bsy$ and provide the associated constants. 

%The assumptions under which this result is shown consist of Gevrey-regularity of the random domain mapping and hit
%associated with the random domain stability analysis Problem \ref{problem:transformed_linearbuckling}. Moreover, we provide the Gevrey-regularity constants induced by the Gevrey-regularity of the domain mapping $\Phi$ in \eqref{eq:definitionPhiandJ}. Note that the random coefficient in $b_\bsy(\bsu,\bsv)$ in \eqref{eq:bilinearforms2} is obtained from solving another source problem of the form $a_{\bsy}(\bsu,\bsv) = \calF_\bsy(\bsv)$. This coupling leads to large constants and thus potential large variation of the operator $b_\bsy$, $\calT_{\bsy}$ and the largest eigenvalue $\mu_{1}(\bsy)$.

Let $\Phi(\bsy)$, $J(\bsy)$ and the space $V$ be defined as in \eqref{eq:definitionPhiandJ} and in \eqref{eq:def_Vspace}, respectively.
%Let $\Phi(\bsy)\colon\Omega\to D(\bsy)$ map the reference domain $\Omega$ to the
%perturbed domain $D(\bsy)=\Phi(\Omega;\bsy)$, and let $J(\cdot;\bsy)=\D_\xi\Phi(\cdot;\bsy)$ be its Jacobian. Moreover, we assume that $\Phi$ does not contain rigid-body motions, that is, some part of the domain with non-zero area remains unchanged under $\Phi$. 
A large class of random domain perturbations is covered by the following assumptions:
\begin{assumption}\label{assumption:GevreyRegularityJ}
\begin{enumerate}
  \item[\textbf{(A1)}]  \textbf{Gevrey regularity of $\Phi$.}
  There are $\delta\ge1$, $C_\Phi>0$ and a sequence $\bsb\in\R_{\ge0}^\infty$ with
  $\bsb\in\ell^p(\N)$ for some $p\in(0,1]$, such that
  \begin{equation}\label{eq:PhiAssumption}  
  \big\|\Dnuy J(\bsy)\big\|_{L^{\infty}(\Omega;\mathbb{T}_2)}
    \le C_\Phi\,\bsb^{\bsnu}\,\fallennudelta
    \qquad\text{for all }\bsy\in\calU,\ \bsnu\in\mathcal F .
  \end{equation}
 \item[\textbf{(A2)}]     \textbf{Uniformly invertible Jacobian.}
  $J(\cdot;\bsy)$ is invertible and the inverse is bounded uniformly as
  \begin{equation}\label{eq:Jacobian_inverseBound}
    \big\|J^{-1}(\bsy)\big\|_{L^\infty(\Omega;\mathbb{T}_2)}\le C_\Phi
    \qquad\text{for all }\bsy\in\calU .
  \end{equation}
  \end{enumerate}
\end{assumption}


\begin{lemma}[Uniform Korn inequality.]\label{lemma:kornconstant}
There is $C_K>0$ such that \tododuc{write on perturbed}
  \begin{equation}\label{eq:uniformKornconstant}
    \|\bsu\|_{H^1(D(\bsy);\R^3)}^2\le C_K\int_{D(\bsy)}\|e(\bsu)\|_F^2\dx
    \qquad\text{for all }\bsy\in\calU,\ \bsu\in H^1(D(\bsy);\R^3),
  \end{equation}
  where $e_{ij}(\bsu) = 1/2 ( \D u_{i}/\D x_j + \D u_{j}/\D x_i)$ is the symmetric gradient on $D(\bsy)$.
  \end{lemma}
  The proof is given in the Appendix~\ref{sec:appendix}.
 % \emph{In words:} the strain energy controls the full $H^1$ norm, with a constant that does not
\begin{example}[Geometric imperfections of the mid-surface]
In the buckling analysis of thin-walled structures, the sensitivity propagation of mid-surface imperfections is of particular interest. For this, let $\omega \subset \R^2 $ be the parameter domain and $ t: \omega \to \R_{>0}$ be the thickness. The mid-surface domain mapping induces a domain mapping of the form
  %A typical example is an affine parametrization of geometric imperfections of the mid-surface,
  \begin{align*}
    \phi(\uxi;\bsy)&=\phi_0(\uxi)+\sum_{i=1}^\infty y_i\,\psi_i(\uxi),
    \\
    \Phi\big((\uxi,\xi_3);\bsy\big)&:=\phi(\uxi;\bsy)+\xi_3\,\vec n(\uxi;\bsy), \qquad  \uxi \in \omega, \vert \xi_3 \vert \leq t(\uxi)/2,   \bsy \in \calU,
  \end{align*}
  where $\vec n(\cdot;\bsy)$ is the unit normal vector of the tangent vectors $ \boldsymbol{a}_{\alpha} :=\D\phi/\D \xi_\alpha$. 
  If the angle of the tangent vectors are uniformly bounded away from zero and the modes satisfy $\sum_i\|\psi_i\|_{W^{2,\infty}(\omega;\R^3)}<\infty$, then
  \eqref{eq:PhiAssumption} holds with weights $b_i\sim\|\psi_i\|_{W^{2,\infty}(\omega;\R^3)}$,
  see \cite{GIANI2025118081}. Here, $W^{2,\infty}(\omega;\R^3)$ denotes the second order infinity Sobolev-norm. Moreover, $C_\Phi$ depends on the largest and smallest eigenvalue of the first, second and third fundamental form of $\phi$, see e.g. \cite[Chapter 2]{ChapelleBathe2011}. This value is dependent on the maximum and minimum extension or contraction and the curvature induced by the parametrization. In fact, one can show that $ \| J \|_2 \sim  \| g_{\alpha \beta } \|_2   $  where $g_{\alpha \beta } = \boldsymbol{g}_\alpha  \cdot \boldsymbol{g}_\beta $ and $\boldsymbol{g}_i = \D\Phi / \D \xi_i$. Here, the metric tensor $g_{\alpha\beta}$ depends on terms of the form $\boldsymbol{a}_{\alpha} \cdot \boldsymbol{a}_{\beta}$, $ \boldsymbol{a}_{\alpha} \cdot \boldsymbol{a}_{\beta,\lambda} $ and $\boldsymbol{a}_{\alpha,\beta} \cdot \boldsymbol{a}_{\lambda,\mu}$.
%  The $W^{2,\infty}$ norm appears because $\vec n$ involves first derivatives of $\phi$,
 % and $\Phi$ is measured in $W^{1,\infty}$.
  %It is needed because otherwise the bilinear form $a_\bsy$ in \eqref{eq:bilinearforms}
  %is not coercive uniformly in $\bsy$.
%\textbf{(A3)} (uniform bounds on the stiffness tensor) gives uniform
%  coercivity and boundedness of $a_\bsy$ in the reference configuration.
%\textbf{(A4)} (uniform Korn constant) makes the coercivity constant of $a_{\bsy}$
%  independent of the shape $D(\bsy)$, and hence of $\bsy$. The coercivity constant plays a key role in the Gevrey-regularity of the operators $\calA_\bsy$ and $m_\bsy$, which is associated with the solution of the prestress $\bssigma_*$.
\end{example}

In this section, we show the following main result, that is, the Gevrey regularity of $\calT_\bsy=\calA_\bsy^{-1}\calM_\bsy$ where $\calA_\bsy$ and $\calM_\bsy$ are defined as in \eqref{eq:DefAy}
and \eqref{eq:DefMy} and $a_\bsy $ and $b_\bsy$ are now the bilinear forms from the buckling problem in \eqref{eq:bilinearforms} and \eqref{eq:bilinearforms2} and $\calF_{\bsy}$ denotes the linear form in \eqref{eq:sourceproblem_weak_star}.
%of the two building blocks $\calA_\bsy^{-1}$ and $\calM_\bsy$ can be proven
%(Lemmas~\ref{lemma:productInversionrule} and~\ref{lemma:coefficientregularitytoOperatorRegularity})%
%carries over to their product $\calT_\bsy=\calA_\bsy^{-1}\calM_\bsy$.
%This is the main result of this section.

\begin{lemma}[Gevrey regularity of $\calT_{\bsy}$]\label{lemma:GevreyRegularityTy}
Let Assumptions \ref{assumption:GevreyRegularityJ} hold. Then $\calT_\bsy=\calA_\bsy^{-1}\calM_\bsy$ is Gevrey regular:
for all $\bsy\in\calU$ and all $\bsnu\in\calF$,
\begin{equation}\label{eq:GevreyTy}
   \|\Dnuy\calT_\bsy\|_{\calL(V,V)}
   \le C_\calT\,\rho_\calT^{|\bsnu|}\,\bsb^{\bsnu}\,\fallennudelta ,
\end{equation}
with constants
\begin{equation}\label{eq:constantRegTy}
    C_\calT=4\,C_{\calA^{-1}}C_\calM,\qquad \rho_\calT=\rho_{\calA^{-1}},
\end{equation}
where $C_{\calA^{-1}}$, $\rho_{\calA^{-1}}$ and $C_\calM$ are defined in
\eqref{eq:constantsInverseAyReg} and \eqref{eq:constantsMyReg}.
 Then $\calT_\bsy$ has bounded variation in the sense of
\eqref{eq:boundvariationTy}, with
\begin{equation*}
    C_V=\,r C_\calT\,\rho_\calT\,\|\bsb\|_{\ell^1(\N)}
\end{equation*}
by Lemma~\ref{lemma:variationboundTy}.
\end{lemma}
\begin{proof}
Using the composition rule for linear operators from Lemma~\ref{lemma:productInversionrule}\,(i) on $\calA^{-1}_\bsy$ and $\calM_{\bsy}$ shows the claim immediately. The main contribution is to show the regularity of the latter terms in Lemma~\ref{lemma:GevreyRegInverseOfAyandFy} and Theorem~\ref{theorem:GevreyRegMy}.
\end{proof}
The constants can be simplified to
\begin{align*}
    C_\calT = \frac{2^{17}}{9} C^3_{\tilde \calC} C_K^2 C_T C_\Phi^5, \quad \rho_{\calT} = 3^9 \cdot 2^4 C_K^2 C_{\tilde \calC}^2 C_\Phi^4,
\end{align*}
where $C_{\tilde \calC} = 2^{11} C_\Phi^5 C_{\calC}$ and $C_\calC$ is the material stiffness tensor bound in \eqref{eq:stiffnesstensorupperbound}. 
%We explain briefly how the constants are derived. 
%We apply the product rule from Lemma~\ref{lemma:productInversionrule}\,(i) with $M=\calA^{-1}$ and $K=\calM$.
%This gives the factor $4\,C_{\calA^{-1}}C_\calM$, and the growth rate $\rho_{\calA^{-1}}^{|\bsnu|}$
%is inherited from the inversion bound (ii).
%The variation bound follows Lemma \ref{lemma:variationboundTy}.
%\end{proof}



\subsection{Gevrey-Regularity of the Bilinear Form Representations} \label{subsec:RegularityA_yM_y} 
In the following section, we assume that Assumption~\ref{assumption:GevreyRegularityJ} and the basic material assumptions in Section~\ref{sec:LinearModel} hold.

We outline the regularity proofs of $\calA^{-1}_{\bsy}$ and $\calM_\bsy$ and show:
\begin{enumerate}
    \item[\textbf{1.}] the regularity of $\tilde \calC(\bsy)$ in \eqref{eq:transformed_stiffness} and $\tilde T^i$ in \eqref{eq:transformedtractionload} in Lemma~\ref{lemma:coefficientGevreyRegularity}.
    \item[\textbf{2.}]  the uniform strong coercivity of $a_{\bsy}(\bsu,\bsv)$ in Lemma~\ref{lemma:coercivityofay} and subsequently the regularity of $\calA^{-1}_\bsy$ in Lemma~\ref{lemma:GevreyRegInverseOfAyandFy}.
    \item[\textbf{3.}]  the regularity of $\bssigma_*$ in Lemma~\ref{lemma:GevreyRegularityPrestress} by using the regularity of
    the solution of the source problem $\bsu^*(\bsy) = \calA_\bsy^{-1}\calF_\bsy$.
    \item[\textbf{4.}] the regularity of $\calM_{\bsy}$ in Theorem~\ref{theorem:GevreyRegMy}.
\end{enumerate}
The regularity of the coefficients in Step \textbf{1} and \textbf{2} are well-known. However, we show the regularity of $\calA_{\bsy}^{-1}$ with a new technique: using Remark~\ref{remark:RegularityCoefficientToBilinearForm} and Lemma~\ref{lemma:RepresenterBilinearFormRegularity}, we obtain the regularity of $\calA_\bsy$. Combined with the uniform strong coercivity of $a_{\bsy}$, we can use Lemma~\ref{lemma:productInversionrule} (ii) to show the regularity of $\calA_{\bsy}^{-1}$.
Step \textbf{3} involves the formula
\begin{align}\label{eq:prestressId2}
    \bssigma_*^{ij}(\bsy) = \calC^{ijks}g_s^l(\bsy) u^*_{k,l}(\bsy), \quad \bsy \in \calU
\end{align}
for the solution of the pre-stress equation. %\eqref{eq:sourceproblem_weak_star}.
Since $\bsu^*(\bsy) = \calA_\bsy^{-1} \calF_{\bsy}$ is a solution to the source problem, this choice satisfies the pre-stress condition in \eqref{eq:sourceproblem_weak_pulled_back}.
Step \textbf{4} then follows with the regularity of $\tilde{\bssigma}$ in \eqref{eq:transforme_prestress} using the product rule in Lemma~\ref{lemma:productInversionrule}(i) and again arguing that the coefficient regularity implies the operator regularity with Remark~\ref{remark:RegularityCoefficientToBilinearForm} and Lemma~\ref{lemma:RepresenterBilinearFormRegularity}.

Coefficient regularity for random domain transforms is a well-known result. To adapt to the notation used in here, we state and prove the results.
\begin{lemma}[Gevrey-regularity of the coefficients of $a_{\bsy},
\calF_{\bsy}$]\label{lemma:coefficientGevreyRegularity}
      For all $\bsnu \in \calF$, we have
     \begin{align*}
         \| \Dnuy \tilde{\calC}^{ijkl} \|_{L^\infty(\Omega;\mathbb{T}_4)}  &\leq  C_{\tilde{\calC}} \rho_{\tilde{\calC}}^{\vert \bsnu \vert}\bsb^{\bsnu} \fallennudelta \\
      %\| \Dnuy \tilde F_{i} \|_{L^\infty(\Omega;\R^{3})}  &\leq  C_{\tilde F_i} \bsb^{\bsnu} \fallennudelta \\
       \| \Dnuy \tilde T^{i} \|_{L^\infty(\partial \Omega;\R^{3})}  &\leq  C_{\tilde T} \rho_{\tilde{T}}^{\vert \bsnu \vert } \bsb^{\bsnu} \fallennudelta
     \end{align*}
     where $\tilde{\calC}$ and $\tilde{T}^i$ are defined in \eqref{eq:transformed_stiffness} and \eqref{eq:transformedtractionload} and the constants are  \begin{align}\label{eq:constantRegTildeC}
         C_{\tilde{\calC}}&= 2^{11} C_{\Phi}^5 C_{\calC}, &  \rho_{\tilde{\calC}} &= 3\cdot 2^4 C^4_\Phi,    \\
     %   \label{eq:constantRegTildeFi}
     %     C_{\tilde F_{i}}&=3\cdot2^7 C_{\Phi}^3 C_{F_{i}}, \\
          \label{eq:constantRegTildeTi}
          C_{\tilde T}&=  2^9  C_{\Phi}^4 C_{T}, &  \rho_{\tilde{T}} &= 3C^2_\Phi.
     \end{align}
\end{lemma}
For the proof, see Appendix \ref{proof:coefficientRegularity}.
The coefficient regularity implies the regularity of $\calA_\bsy$ and $\calM_\bsy$ with the \emph{same} constants.
\begin{remark}\label{remark:RegularityCoefficientToBilinearForm}
Let $\tilde{\calC}^{ijkl}(\bsy)$ as in \eqref{eq:transformed_stiffness} and $\tilde{\sigma}_*^{ij}(\bsy)$ as in \eqref{eq:transforme_prestress} be Gevrey-regular, that is, 
there exist $C_1,C_2 \geq 0$ and $\bsb_1, \bsb_2 \in \R^\infty_{\geq 0}$ such that
\begin{align}\label{eq:Gevrey_assump1}
     \|\D^{\bsnu} \tilde{\calC}^{ijkl}(\bsy) \|_{L^\infty(\Omega;\mathbb{T}_4)} &\leq C_1 \bsb_1^\bsnu \fallennudelta,  \\
     \label{eq:Gevrey_assump2}
      \| \D^{\bsnu}\tilde{\bssigma}_*^{ij}(\bsy) \|_{L^\infty(\Omega;\mathbb{T}_2)} &\leq  C_2 \bsb_2^\bsnu \fallennudelta, 
\end{align}
for all $\bsnu \in \calF, \bsy \in \calU$.
Then the bilinear forms from $  a_\bsy(\bsu, \bsv) $in \eqref{eq:bilinearforms} and  $ b_\bsy(\bsu, \bsv)  $  in \eqref{eq:bilinearforms2} for $\bsu, \bsv \in V$ are Gevrey-regular with the same constants, that is,
\begin{align*}
     \|\D^{\bsnu}a_\bsy \|_{\calL(V \times V) } &\leq C_1 \bsb_1^\bsnu \fallennudelta,  \\
      \|\D^{\bsnu} b_\bsy \|_{\calL(V \times V) } &\leq  C_2 \bsb_2^\bsnu \fallennudelta, 
\end{align*} for all $\bsnu \in \calF, \bsy \in \calU$.
\end{remark}
\begin{proof}
    The norms on $\mathbb{T}_4$ and on $\mathbb{T}_2$ are
    \begin{align*}
        \| M^{ijkl}\|_{\mathbb{T}_4} = \sup_{ \| A \|_F  = \| B \|_F = 1}    M^{ijkl} A_{ij}B_{kl}, \quad \text{ and } \quad 
        \| M^{ij} \|_{\mathbb{T}_2} = \sup_{ \|u \|_2 = \| v \|_2 =1} M^{ij}v_i v_j,
    \end{align*}
    respectively, where $F$ denotes the Frobenius norm on matrices and $\| \cdot \|_2$ is the Euclidean norm. 
    Since the derivatives of $\tilde{\bssigma}^{ij}_*(\bsy)$ and $\tilde{\calC}^{ijkl}(\bsy)$ are bounded by \eqref{eq:Gevrey_assump1} and \eqref{eq:Gevrey_assump2}, we can interchange differentiation and integration and obtain
    \begin{align*}
         \| \D^{\bsnu } a_\bsy \|_{\calL(V \times V) } &=  \sup_{ \| \bsu \|_V = \| \bsv \|_V =1  } \D^{\bsnu} \int_{\Omega} \tilde{\calC}^{ijkl}(\bsy) u_{i,j} v_{k,l} \dxi =    \sup_{ \| \bsu \|_V = \| \bsv \|_V =1  } \int_{\Omega} \D^{\bsnu} \tilde{\calC}^{ijkl}(\bsy) u_{i,j} v_{k,l} \dxi
         \\
         &\leq  \| \D^{\bsnu} \tilde{\calC} (\bsy)\|_{L^\infty(\Omega;\mathbb{T}_4) } \sup_{ \| \bsu \|_V = \| \bsv \|_V =1  } \int_{\Omega}  \| u_{i,j} \|_F \|v_{k,l} \|_F \dxi \\
         &\leq   \| \D^{\bsnu} \tilde{\calC} (\bsy)\|_{L^\infty(\Omega;\mathbb{T}_4) },
    \end{align*}
    where the Cauchy-Schwarz inequality was applied to obtain the last inequality.
    The same computation yields $  \| \D^{\bsnu } b_\bsy \|_{\calL(V \times V) }   \leq\| \D^{\bsnu} \tilde{\bssigma}_*^{ij} (\bsy)\|_{L^\infty(\Omega;\mathbb{T}_2) }.$
    %Furthermore, in both computations, we require that the derivatives of $\tilde{\bssigma}^{ij}_*(\bsy)$ and $\tilde{\calC}^{ijkl}(\bsy)$ exist, which is guaranteed by \eqref{eq:Gevrey_assump1} and \eqref{eq:Gevrey_assump2}.
\end{proof}



From Remark~\ref{remark:RegularityCoefficientToBilinearForm} and Lemma~\ref{lemma:RepresenterBilinearFormRegularity}, we immediately obtain the Gevrey-regularity of $\calA_{\bsy}$.
\begin{lemma}[Gevrey-Regularity of $\calA_\bsy$]\label{lemma:GevreyRegularityAy}
    For all $\bsnu \in \calF$, we have 
    \begin{align*}
        \| \D^\bsnu \calA_\bsy \|_{\calL(V,V')}  \leq C_{\calA} \rho_{\calA}^{\vert \bsnu \vert} \bsb^{\bsnu} \fallennudelta,
    \end{align*}
    for all $\bsy \in \calU$ and where
    \begin{align}\label{eq:constantsAyReg}
        C_\calA = C_K C_{\tilde \calC} ,\quad \rho_\calA = \rho_{\tilde \calC}
    \end{align}
    where $C_K$ is the uniform Korn constant from \eqref{eq:uniformKornconstant}.
\end{lemma}
Due to Lemma~\ref{lemma:productInversionrule}(ii) we need $C_{\calA}$ to bound $\| \calA_{\bsy}^{-1} \|_{\calL(V',V)}$ uniformly. The above choice satisfies this property:
\begin{lemma}[Uniform strong coercivity of $a_\bsy$]\label{lemma:coercivityofay}
     $a_\bsy$ is uniformly strongly coercive with the inverse of $C_{\calA}$ defined in \eqref{eq:constantsAyReg}, that is, for all $\bsy \in \calU$
    \begin{align*}
        a_{\bsy}(\bsu,\bsu) \geq C_{\calA}^{-1} \| \bsu \|_V, \quad \forall  \bsu \in V.
    \end{align*}
\end{lemma}
\begin{proof}
    The proof of coercivity of $\calA_\bsy$ is similar to \cite[Lemma 2.7]{Harbrecht2023DomainUncertainty}. However, the uniform Korn constant can exceed $2$ for general domain shapes and Dirichlet boundary $\D \Omega_D$. We repeat the same argument and adopt the Korn inequality \eqref{eq:uniformKornconstant}. We have
\begin{align*}
    a_{\bsy}(\bsu,\bsu) &= \int_{D(\bsy)} \calC^{ijkl} \DD{ u_i}{ x_j} \DD{u_k}{x_l} \dx \geq C^{-1}_{\calC}  \int_{D(\bsy)}  \|e_{ij}(\bsu)  \|_F^2
    \\
    &\geq C_K^{-1}  C^{-1}_{\calC} \| u \circ \Phi(\cdot;\bsy)\|^2_{H^{1}(D(\bsy);\R^3)} \geq  C_K^{-1}  C^{-1}_{\calC} C_{\Phi}^{-5}  \| u \|^2_{V}  \geq C_\calA^{-1} \| u \|^2_{V},
\end{align*} for all $\bsu \in V$ and $\bsy \in \calU$. The first inequality comes from the material assumption \eqref{eq:stiffnesstensorcoercivity} and the second uses the Korn inequality \eqref{eq:uniformKornconstant}. Note that the third inequality comes from the estimates
\begin{align*}
    \| J^{-\top} (\bsy) \|_{L^\infty(\Omega;\mathbb{T}_2)}^2 \geq C^{-2}_{\Phi} \text{ and }  \| \det J(\bsy)\|_{L^\infty(\Omega)} \geq C_\Phi^{-3}, \quad  \forall \bsy \in \calU.
\end{align*} 
\end{proof}
By the Riesz representation theorem $\calA_{\bsy}^{-1} $ exists for every $\bsy \in \calU$ and is moreover bounded by $\| \calA_{\bsy}^{-1} \|_{\calL(V',V)} \leq C_{\calA}$. Hence the next lemma is a direct consequence of applying Lemma~\ref{lemma:productInversionrule}(ii).
\begin{lemma}[Gevrey-Regularity of $\calA_{\bsy}^{-1}$ and $\calF_{\bsy}$]\label{lemma:GevreyRegInverseOfAyandFy}
    For all $\bsnu \in \calF$ and $\bsy \in \calU$
    \begin{align*}
        \| \calA_{\bsy}^{-1} \|_{\calL(V,V')}& \leq C_{\calA^{-1}} \rho_{\calA^{-1}}^{\vert \bsnu \vert} \bsb^{\bsnu} \fallennudelta, \\
        \| \calF_\bsy \|_{\calL(V,V')} & \leq C_{\calF} \rho_{\calF}^{\vert \bsnu \vert} \bsb^{\bsnu} \fallennudelta,
    \end{align*}
    where
    \begin{align} \label{eq:constantsInverseAyReg}
        C_{\calA^{-1}} &= C_{\calA}/3 ,\quad \rho_{\calA^{-1}} = (3C_{\calA})^2 \rho_{\calA}, \\
        \label{eq:constantsFyReg}
        C_\calF &= C_{\tilde T} ,\quad \rho_\calF = \rho_{\tilde T}.
    \end{align}
\end{lemma}
\begin{proof} We only need to show the statement for $\calF_\bsy$ in a similar way as in Remark~\ref{remark:RegularityCoefficientToBilinearForm}.
    Note that $ \| \D^\bsnu  \calF_{\bsy} \|_{V'} = \sup_{ \|v \|_V = 1 } \D^\bsnu   \oint_{\partial \Omega_T} \tilde T^i(\bsy) v_i \mathop{dS} $. And since $ \tilde T^i(\bsy)$ is Gevrey-regular, we can interchange integration and differentiation and estimate the variation of traction in the $L^\infty$-norm: For any $\bsy \in \calU$ and $\bsnu \in \calF$, we have
    \begin{align*}
        \| \D^\bsnu  \calF_{\bsy} \|_{V'} \leq    \sup_{ \|v \|_V = 1 }  \|\D^\bsnu \tilde T^i(\bsy) \|_{L^\infty(\partial \Omega_T; \R^3)} \oint_{\D\Omega_T} || v_i||_1 \mathop{dS} \leq  \|\D^\bsnu \tilde T^i(\bsy) \|_{L^\infty(\partial \Omega; \R^3)}.
    \end{align*}
\end{proof}
\begin{lemma}[Gevrey-regularity of $\tilde \bssigma$]\label{lemma:GevreyRegularityPrestress}
The transformed prestress $\tilde \bssigma$ defined in \eqref{eq:transforme_prestress} where $\bssigma_*$ has the identity \eqref{eq:prestressId2} is Gevrey-regular,
that is, for all $\bsnu \in \calF$ and $\bsy \in \calU$
    \begin{align*}
        \| \tilde  \bssigma  \|_{L^\infty(\Omega;\mathbb{T}_2)} & \leq C_{\tilde  \bssigma} \rho_{\tilde  \bssigma}^{\vert \bsnu \vert} \bsb^{\bsnu} \fallennudelta,
    \end{align*}
    where
    \begin{align} 
        \label{eq:constantsTransformedPrestressReg}
        C_{\tilde  \bssigma}= 2^6 C_{\tilde \calC} C_{\Phi} C_{\calA^{-1}} C_{\calF},\quad \rho_{\tilde  \bssigma} = \rho_{\calA^{-1}}.
    \end{align}
\end{lemma}
\begin{proof}
The relevant identities are in \eqref{eq:prestressId2} and \eqref{eq:transforme_prestress}:
\begin{align*}
        \bssigma_*^{ij}(\bsy) &= \calC^{ijks}g_s^l(\bsy) u^*_{k,l}(\bsy), \\
        \tilde \bssigma^{ij}(\bsy) &= \sqrt{g(\bsy)} \bssigma_*^{rs}(\bsy) g_{r}^i(\bsy)  g_{s}^j(\bsy) 
\end{align*}
where $\bsu^* = \calA_{\bsy}^{-1}\calF_\bsy$.
Observe that $\bsu^*(\bsy)$ is Gevrey-regular with constants $C_{\bsu^*} = 4 C_{\calA^{-1}} C_{\calF}$ and $\rho_{\bsu^*} = \rho_{\calA^{-1}}$.
The claim follows by using the product rule in Lemma~\ref{lemma:productInversionrule}(i) two more times on
\begin{align*}
    \tilde \bssigma^{ij}(\bsy) &=  \tilde \calC^{ijkr} g_{r}^l(\bsy) \bsu^*_{k,l}(\bsy),
\end{align*}
and by observing that the product rule chooses $\rho$ according to the maximum $\rho$ value of the single terms in the product.
\end{proof}
Note that by Remark~\ref{remark:RegularityCoefficientToBilinearForm} and Lemma~\ref{lemma:RepresenterBilinearFormRegularity}, the regularity of $\calM_{\bsy}$ is an immediate consequence.
\begin{theorem}[Gevrey-regularity of $\calM_\bsy$]\label{theorem:GevreyRegMy}
     For all $\bsnu \in \calF$, we have 
    \begin{align*}
        \| \D^\bsnu \calM_\bsy \|_{\calL(V,V')}  \leq C_{\calM} \rho_{\calM}^{\vert \bsnu \vert} \bsb^{\bsnu} \fallennudelta,
    \end{align*}
    for all $\bsy \in \calU$ and where
    \begin{align}\label{eq:constantsMyReg}
       C_{\calM} = C_{\tilde \bssigma} ,\quad\rho_{\calM} = \rho_{\tilde  \bssigma}.
    \end{align}
\end{theorem}
    




\section{Numerical Example}
\label{sec:numerical-example}
% Measured data: application-vshp2-eigen-buckling/experiments/output/
% bunter-first-set/theory-evidence.json; archive and replay identities retained there.
\subsection{One-dimensional profile}
\begin{gather}
 y=(y_1,y_2)\in[-1,1]^2,\qquad
 \mathrm{d}\pi(y)=\tfrac14\,\mathrm{d}y_1\mathrm{d}y_2,
 \qquad s\in[-L,L],\quad L=2,\quad b=\tfrac7{20},\\
 \psi_1(s)=\begin{cases}(1-(s/b)^2)^4,&|s|<b,\\0,&|s|\ge b,\end{cases}
 \qquad
 \psi_2(s)=\frac{19(s/b)^2-1}{c_*}\psi_1(s),
 \quad c_*=\frac{483729408}{407253125},\\
 \psi_i\in C^3([-L,L]),\qquad
 \|\psi_i\|_{L^\infty}=1,\qquad
 \int_{-L}^{L}\psi_1(s)\psi_2(s)\,\mathrm{d}s=0,\\
 f(s;y)=1-\frac1{20}\psi_1(s)
       +\frac{y_1}{100}\psi_1(s)+\frac{y_2}{400}\psi_2(s),
 \qquad \partial_\varphi f=0.
 \label{eq:experiment-profile}
\end{gather}
\begin{gather}
 X=(s,f\cos\varphi,f\sin\varphi),\qquad
 a=\sqrt{1+f_s^2},\qquad
 e_1=a^{-1}(1,f_s\cos\varphi,f_s\sin\varphi),\\
 e_2=(0,-\sin\varphi,\cos\varphi),\qquad
 n=a^{-1}(-f_s,\cos\varphi,\sin\varphi),\\
 \Gamma_y=X([-L,L]\times[0,2\pi)),\qquad Y=X+zn,\\ |z|<t/2,\quad t=\tfrac1{100},\qquad
 \mathrm{d}A=af\,\mathrm{d}s\,\mathrm{d}\varphi,\\
 D_1=a^{-1}\partial_s,\quad D_2=f^{-1}\partial_\varphi,\quad
 c_1=-f_{ss}/a^3,\quad c_2=(af)^{-1},\quad \gamma=f_s/(af),\\
 \lambda_{\mathrm L}=\mu=1,\qquad E=\tfrac52,\qquad\nu=\tfrac14.
\end{gather}

\subsection{Five-field model and prestress}
\begin{gather}
 U=(u,v,w,\theta,\psi),\quad
 q=u e_1+v e_2+w n,\quad \eta=-\theta e_1-\psi e_2,
 \quad\mathcal U(Y)=q+z\eta,\\
 \begin{aligned}
 \beta_{11}&=D_1u+c_1w,&\beta_{22}&=D_2v+\gamma u+c_2w,\\
 2\beta_{12}&=D_1v+D_2u-\gamma v,\\
 \kappa_{11}&=D_1\theta,&\kappa_{22}&=D_2\psi+\gamma\theta,\\
 2\kappa_{12}&=D_1\psi+D_2\theta-\gamma\psi
                 -c_1(D_2u-\gamma v)-c_2D_1v,\\
 \rho_1&=D_1w-c_1u-\theta,&\rho_2&=D_2w-c_2v-\psi.
 \end{aligned}
\end{gather}
\begin{gather}
 \mathbb C_{\mathrm{ps}}A=\frac{E}{1+\nu}
       \left(A+\frac{\nu}{1-\nu}\operatorname{tr}(A)I_2\right),\\
 a_y^{\mathrm N}(U,V)=\int_{\Gamma_y}\left[
 t\mathbb C_{\mathrm{ps}}\beta(U):\beta(V)
 +\frac{t^3}{12}\mathbb C_{\mathrm{ps}}\kappa(U):\kappa(V)
 +\mu t\rho(U)\cdot\rho(V)\right]\mathrm{d}A.
 \label{eq:experiment-naghdi-form}
\end{gather}
\begin{gather}
 U_h^D=(u_h^D,0,w_h^D,\theta_h^D,0),\qquad
 u_h^D(\pm L)=\mp L/200,\quad\theta_h^D(\pm L)=0,\\
 a_y^{\mathrm N}(U_h^D,V_h)=0
 \quad\bigl(V_h=(v_1,0,v_3,v_4,0),\ v_1(\pm L)=v_4(\pm L)=0\bigr),\\
 R_h^D=\frac{\partial}{\partial u_h^D(L)}
                  \frac12a_y^{\mathrm N}(U_h^D,U_h^D),\qquad
 F_0=2\pi t=\frac\pi{50},\qquad c_h=\frac{F_0}{|R_h^D|},\\
 N^0=c_h t\mathbb C_{\mathrm{ps}}\beta(U_h^D),\qquad
 B^0=c_h\frac{t^3}{12}\mathbb C_{\mathrm{ps}}\kappa(U_h^D),\qquad
 Q^0=c_h\mu t\rho(U_h^D),\\
 \widehat\sigma^0_{\alpha\beta}(z)=N^0_{\alpha\beta}/t
                       -12zB^0_{\alpha\beta}/t^3,\quad
 \widehat\sigma^0_{\alpha3}=\widehat\sigma^0_{3\alpha}=Q^0_\alpha/t,
 \quad\widehat\sigma^0_{33}=0,\quad \alpha,\beta\in\{1,2\},\\
 \int_{-t/2}^{t/2}\widehat\sigma^0_{\alpha\beta}\,\mathrm{d}z=N^0_{\alpha\beta},
 \quad\int_{-t/2}^{t/2}z\widehat\sigma^0_{\alpha\beta}\,\mathrm{d}z=-B^0_{\alpha\beta},
 \quad\int_{-t/2}^{t/2}z^2\widehat\sigma^0_{\alpha\beta}\,\mathrm{d}z
       =\frac{t^2}{12}N^0_{\alpha\beta}.
 \label{eq:experiment-stress}
\end{gather}
\begin{gather}
 H_{i\alpha}(U,z)=e_i\cdot(D_\alpha q+zD_\alpha\eta),\quad
 H_{i3}(U,z)=e_i\cdot\eta,\quad(e_3=n),\\
 g_y^{\mathrm N}(U,V)=-\int_{\Gamma_y}\int_{-t/2}^{t/2}
 \sum_{i,j,k=1}^3\widehat\sigma^0_{jk}(z)H_{ij}(U,z)H_{ik}(V,z)
 \,\mathrm{d}z\,\mathrm{d}A,\\
 a_y^{\mathrm N}(\Phi_{j,h},V_h)=\lambda_{j,h}(y)g_y^{\mathrm N}(\Phi_{j,h},V_h),
 \qquad0<\lambda_{1,h}\le\lambda_{2,h}\le\cdots .
 \label{eq:experiment-pencil}
\end{gather}

\subsection{Discrete space and sampled subspace}
\begin{gather}
 \mathcal T_s:\quad[-L,-b],[-b,b],[b,L],\qquad
 (n_-,n_0,n_+)=(12,12,12),\qquad p=5,\\
 S_h^5=\{v\in C^0([-L,L]):v|_I\in\mathbb P_5(I),\ I\in\mathcal T_s\},\\
 \mathcal F_{24}=\operatorname{span}\{1,\cos k\varphi,\sin k\varphi:1\le k\le24\},\\
 \mathcal R=\operatorname{span}\{(0,1,0),(0,0,1),(1,0,0)\mathbin\times Y\},\\
 \mathcal W_h=\{\mathcal U(U_h):U_h\in(S_h^5\otimes\mathcal F_{24})^5\}+\mathcal R,\\
 P_A=Y(-L,0,t/2),\quad P_B=Y(-L,\pi/2,t/2),\\
 C_h\mathcal V_h=
 \bigl(\mathcal V_{h,2}(P_A),\mathcal V_{h,3}(P_A),\mathcal V_{h,3}(P_B)\bigr)^T=0,
 \quad\mathcal V_{h,1}|_{s=\pm L}=0,\\
 \mathcal E_{4,h}(y)=\operatorname{span}\{\Phi_{1,h}(y),\ldots,\Phi_{4,h}(y)\}.
\end{gather}
\begin{gather}
 P_q(x)=\frac{1}{2^q q!}\frac{\mathrm d^q}{\mathrm dx^q}(x^2-1)^q,\qquad
 \mathcal G_q=\{x\in(-1,1):P_q(x)=0\},\\
 \mathcal B=\{-1,0,1\}^2\setminus\{(0,0)\},\qquad
 \mathcal S=\mathcal B\cup\bigcup_{q\in\{3,5,7,9\}}\mathcal G_q^{\otimes2},\\
 \mathcal L_{m,h}(y)=\{\!\{(k_j,r_j)_{\mathrm{dominant}}:1\le j\le m\}\!\},\\
 \sum_{q\in\{3,5,7,9\}}q^2=164,\quad |\mathcal B|=8,\\
 \mathcal L_{4,h}(y)=\{\!\{(3,1),(3,1),(4,1),(4,1)\}\!\}\quad(y\in\mathcal S),\\
 g_{4,h}(y)=\lambda_{5,h}(y)-\lambda_{4,h}(y),\qquad
 \widehat g_{4,h}(y)=g_{4,h}(y)/\lambda_{4,h}(y).
\end{gather}
\begin{table}[htbp]
\centering\small
\caption{Sampled spectra; dominant Fourier labels $(k,r)$; $r_{j,h}=\|K_h\phi_{j,h}-\lambda_{j,h}G_h\phi_{j,h}\|_2/(\|K_h\phi_{j,h}\|_2+|\lambda_{j,h}|\|G_h\phi_{j,h}\|_2)$.}\label{tab:theory-gaps}
\begin{tabular}{lcc}
\hline
Statistic & Value & $y$\\
\hline
$\min_{y\in\mathcal S}g_{4,h}(y)$ & $6.8230968\times10^{-4}$ & $(-1,-1)$\\
$\min_{y\in\mathcal S}\widehat g_{4,h}(y)$ & $0.110556687$ & $(-1,0)$\\
$\mathcal L_{6,h}(y)=\mathcal L_{6,h}(0,0)$ & $\mathrm{false}$ & $(1,-1)$\\
$\max_{q\in\{3,5,7,9\},\ y\in\mathcal G_q^{\otimes2},\ j\leq24} r_{j,h}(y)$ & $1.06248\times10^{-11}$ & ---\\
\hline
\end{tabular}
\end{table}
\subsection{Ordered response statistics}
\begin{gather}
 \{(y_\ell^{(q)},w_\ell^{(q)})\}_{\ell=1}^{q^2},\qquad
 y_\ell^{(q)}\in\mathcal G_q^{\otimes2},\quad
 w_\ell^{(q)}>0,\quad\sum_\ell w_\ell^{(q)}=1,\\
 m_j^{(q)}=\sum_\ell w_\ell^{(q)}\lambda_{j,h}(y_\ell^{(q)}),\qquad
 v_j^{(q)}=\sum_\ell w_\ell^{(q)}
       (\lambda_{j,h}(y_\ell^{(q)})-m_j^{(q)})^2,\\
 \delta_{m,j}=\frac{|m_j^{(9)}-m_j^{(7)}|}{|m_j^{(7)}|},\qquad
 \delta_{v,j}=\frac{|v_j^{(9)}-v_j^{(7)}|}{v_j^{(7)}}.
\end{gather}
\begin{table}[htbp]
\centering\small
\caption{Tensor Gauss--Legendre statistics; $(p,n_s,k_{\max})=(5,36,24)$.}\label{tab:theory-moments}
\begin{tabular}{rrrrr}
\hline
$j$ & $10^3m_j^{(9)}$ & $10^4\sqrt{v_j^{(9)}}$ & $\delta_{m,j}$ & $\delta_{v,j}$\\
\hline
1 & 4.94051103 & 6.8096656 & $9.672\times10^{-6}$ & $1.964\times10^{-4}$\\
2 & 4.94086278 & 6.8075369 & $6.204\times10^{-6}$ & $2.016\times10^{-4}$\\
3 & 5.16663021 & 7.2064533 & $9.249\times10^{-6}$ & $2.171\times10^{-4}$\\
4 & 5.16881603 & 7.2233638 & $5.931\times10^{-6}$ & $2.059\times10^{-4}$\\
5 & 6.16087825 & 5.6798026 & $6.168\times10^{-13}$ & $2.046\times10^{-11}$\\
6 & 6.20215706 & 5.7218973 & $1.287\times10^{-7}$ & $4.655\times10^{-6}$\\
\hline
\end{tabular}
\end{table}

\subsection{Truncated Gaussian distribution(fixed)}
We now consider the kernel function $f_{\sigma}(x)=\exp\brac{-\frac{x^2}{2\sigma}}$. For a fixed $c > 0$ and the corresponding truncated domain $\Omega=[-c,c]$, the probability density function (pdf) and the cumulative distribution function (cdf) are defined as follows:
\begin{align}
    \text{pdf: }& \phi(x) := \frac{f_{\sigma}(x)}{\int_{\Omega}f_{\sigma}(s) \,ds}, \\
    \text{cdf: }& \Phi(x) := \int_{-c}^x \phi(s) \,ds.
\end{align}
Since $\Phi: [-c,c] \to [0,1]$ is a bijection, its inverse $\Phi^{-1}: [0,1] \to [-c,c]$ is well-defined. We define the function $g_1: [-\tfrac{1}{2},\tfrac{1}{2}] \to [-c,c]$ as the shifted inverse of the cdf, given by:
$$
g_1(y) = \Phi^{-1}(y+\tfrac{1}{2}).
$$
Because the kernel function $f_{\sigma}$ is analytic, it follows that the pdf $\phi$, the cdf $\Phi$, and the shifted inverse $g_1$ are all analytic. We can consider a different kernel, such as $f_2(x) = 1 - \exp\brac{-\tfrac{1}{x^2}}$, to construct a Gevrey class function $g_2 \in G^{\delta}$ ($\delta=\tfrac{3}{2}$) following the same approach.
Then, we could modify the function \ref{eq:experiment-profile} as follows
\begin{align}
    f^{(j)}(s,y)= 1-\frac1{20}\psi_1(s)
       +v\left(g_j(y_1)\psi_1(s)+\frac{1}{4}g_j(y_2)\psi_2(s) \right).
\end{align}
Consider the previous case for $v = 1/100$. Then
varying $c >0$ and $v >0$ should exhibit the regime in which mixing of the smallest eigenvalue cluster occurs. 

\section{Discussion}\label{sec:Discussion}
\subsection*{Conclusion}
\subsection*{Acknowledgments}
\section{Appendix}\label{sec:appendix}
\begin{proof}[Proof of Lemma \ref{lemma:coefficientGevreyRegularity}]\label{proof:coefficientRegularity}
    From the Gevrey-regularity of $\Phi$ in Assumption \eqref{eq:PhiAssumption}, we have
    \begin{align*}
        \| \Dnuy  J \|_{L^\infty (\Omega;
        \mathbb{M}_2^3)} \leq C_{\Phi} \bsb^{\bsnu} \fallennudelta 
    \end{align*}
    and with the product rule, inversion, and the determinant in Lemma \ref{lemma:productInversionrule} or \cite[Lemma 2]{chernov2026gevreyclassregularitysteadystate},
    we obtain
    \begin{align*}
        \| \Dnuy J^{-1} (\cdot;\bsy) \|_{L^\infty (\Omega;\mathbb{M}_2^3)}  &\leq \frac{C_\Phi}3  (3 C_\Phi^2)^{\vert \bsnu \vert } \bsb^{\bsnu} \fallennudelta,   \\
        \| \Dnuy  (J^\top J) (\cdot;\bsy) \|_{L^\infty (\Omega;\mathbb{M}_2^3)} &\leq 2^2 C_{\Phi}^2 \bsb^{\bsnu} \fallennudelta,  \\
        \| \Dnuy  (J^\top J)^{-1} (\cdot;\bsy)\|_{L^\infty (\Omega;\mathbb{M}_2^3)} &\leq \frac{2^2C_{\Phi}^2}{3} \left(3 \cdot 2^4C_{\Phi}^4 \right)^{\vert \bsnu \vert }\bsb^{\bsnu} \fallennudelta,  \\
        \| \Dnuy  \det J (\cdot;\bsy) \|_{L^\infty (\Omega;\mathbb{M}_2^3)} &\leq 3 \cdot 2^7 C_{\Phi}^3 \bsb^{\bsnu} \fallennudelta  
    \end{align*}
    for all $\bsy \in \calU$.
    Now, recalling the following identities of the transformed coefficients in Section \ref{section:linearbuckingformulation}, we obtain
    \begin{align*}
         \| \Dnuy \tilde{\calC}_{ijkl} \|_{L^\infty(\Omega;\R^4)} &\leq \| \Dnuy \left[ \det J J_{ri}^{-1} J_{js}^{-1} C_{rjsl} \right]  \|_{L^\infty(\Omega;\mathbb{M}_4^3)} \leq  2^{11} C^5_{\Phi} (3\cdot 2^4 C^4_\Phi)^{\vert \bsnu \vert } \bsb^\bsnu \fallennudelta,\\
         % \| \Dnuy \tilde F_{i} \|_{L^\infty(\Omega;\R^{3})} &\leq \| \Dnuy \det J  F_i \|_{L^\infty(\Omega;\R^3)} \leq 3 \cdot 2^7 C^3_\Phi \bsb^\bsnu \fallennudelta, \\ 
          \| \Dnuy \tilde T_{i} \|_{L^\infty(\partial\Omega;\R^{3})} &\leq \| \Dnuy \left[ \det J \| J^{-\top} n_{\text{ref}} \|_2  T_i \right] \|_{L^\infty(\partial\Omega;\R^3)} \leq 2^9 C^4_\Phi (3C^2_\Phi)^{\vert \bsnu \vert} \bsb^\bsnu \fallennudelta,
    \end{align*}
  %  Combining the last two inequalities with the product rule, \eqref{eq:stiffnesstensorupperbound}, \eqref{eq:bodyloadbound} and \eqref{eq:tractionloadbound}, respectively,
  as claimed.
\end{proof}

\begin{proof}[Proof of Lemma~\ref{lemma:kornconstant}]\label{proof:KornConstant} For notational convenience, assume that the function spaces are combined with standard vector and matrix norms, i.e., Euclidean and Frobenius norm, respectively.

Suppose \eqref{eq:uniformKornconstant} is not true.
    Then there exists a sequence $ \{ (\bsu_n,\bsy_n) \} \subseteq V(\calE(\bsy_n)) \times \calU$
    such that
    \begin{align}\label{eq:strainscontradiction}
        \| e(\bsu_n) \|_{L^2(\calE_n)} < \frac 1 n \| \bsu_n \|_{H^1(\calE_n)},
    \end{align}
    where we set $\| \bsu_n \|_{H^1(\calE_\infty)} = 1$ and $\calE_n = \calE(\bsy_n)$ for $n \in \N$. \\
    Since $\calU $ is compact there exists a subsequence of $\bsy_n$ that converges to $\bsy \in \calU$.
    Let $\calE = \calE(\bsy)$ and consider the subsequence in the following.\\
We may extend each $\bsu_n$  to $\calE_\infty$ with the Sobolev extension theorem.
Since $\calE(\bsy)$ is compact for all $\bsy \in \calU$ and $\calU$ is compact, $\calE_\infty$ is compact.

Hence there exists $\bsu \in V(\calE_\infty)$ such that
\begin{align*}
\bsu_n \rightharpoonup \bsu \quad \text{in } H^1(\calE_\infty), \quad \text{and } \| \bsu_n\|_{H^1(\calE_\infty)} \leq D
\end{align*}
for some $D\geq 1$. Again, restrict to such a subsequence.

By Assumption~\ref{assumption:GevreyRegularityJ} $\bsy \mapsto \calE(\bsy)$ is continuous and hence $\vert \calE \Delta \calE_n \vert  \to 0 $ for $n \to \infty$. Thus, by \eqref{eq:strainscontradiction},
\begin{align*}
\| e(\bsu_n) \|^2_{L^2(\calE)} =  \| e(\bsu_n) \|^2_{L^2(\calE \cap \calE_n)}  + \| e(\bsu_n) \|^2_{L^2(\calE \setminus \calE_n)}   \leq  \| e(\bsu_n) \|^2_{L^2(\calE_n)} + \| D\|^2_{L^2(\calE \setminus \calE_n)}  \to  0.
\end{align*}
Hence, $ \| e(\bsu) \|^2_{L^2(\calE)} = 0 $ and the classical Korn inequality \cite[Theorem 1.6-1]{ciarletIntroductionToShellTheory} asserts
\begin{align}\label{eq:uiszero}
\| \bsu \|_{H^1(\calE) } = 0.
\end{align}
However, since
\begin{align*}
 1 = \| \bsu_n\|^2_{H^1( \calE_n)} = \| \bsu_n \|^2_{H^1(\calE \cap \calE_n)} + \| \bsu_n \|^2_{H^1(\calE_n \setminus \calE )}
\end{align*}
and since the second term goes to $0$, there exists an $n_0 $ such that
\begin{align*}
 \| \bsu_n\|_{H^1(\calE \cap \calE_n)}   \geq \frac 1 2, \quad \forall n > n_0
\end{align*}

which leads to a contradiction with \eqref{eq:uiszero} by
\begin{align*}
\|\bsu \|_{H^1(\calE)} = \lim_{n \to \infty } \|\bsu_n \|_{H^1(\calE)} \geq \lim_{n \to \infty } \|\bsu_n \|_{H^1(\calE_n \cap \calE)} \geq \frac 1 2.
\end{align*}
\end{proof}

%\subsection{Asymptotic Analysis of Domain Perturbation Sensitivity of the Critical Buckling Load}

The idea of this section is based on Dauge and Suri in \cite{dauge:hal-00012083}. Here, we generalize their asymptotic analysis for curvilinear coordinates, exhibiting the effect of random domain perturbations when the thickness of the shell goes to zero.
\paragraph{Setting.} For the asymptotic analysis, we define the pullback to the uniform reference domain. Let $\omega \subseteq \R^2 $ be a bounded domain. 
For thickness $\epsilon > 0$, let $\Omega^\epsilon := \omega \times (-\epsilon, \epsilon)$ and let $\Omega = \omega \times (-1,1)$ denote the uniform reference domain.

We write coordinates in $\Omega^\epsilon$ as $\vxi = (\xi^1, \xi^2, \xi^3)^\top$ and in $\Omega$ as $\vec \eta  = (\eta_1,\eta_2,\eta_3)^\top$. Moreover, we use the basic notation for covariant, contra-variant bases, fundamental forms and volume element 
as follows:
\begin{align*}
    \va_{\alpha} &= \DD{\phi}{\xi^\alpha} ,\quad  \va_3 = \va_1 \times \va_2 / \|  \va_1 \times \va_2  \|_2, \quad, \vec g_i  = \DD{\Phi}{\xi^i}, \quad \vec g^{i} \cdot \vec g_{j}= \delta_{ij}, \quad g^{ij}= \vec g^{i} \cdot \vec g^{j} \\\ 
    a_{\alpha\beta} &= \va_\alpha \cdot \va_\beta, \quad \va^\alpha \cdot \va_\beta = \delta_{\alpha\beta},\quad  a^{\lambda \mu}= \va^\lambda \cdot \va^\mu,
    \quad g_{ij} = \vec g_{i} \cdot \vec g_{j}, \quad  b_{\alpha\beta } = \va_3 \cdot \DD {\va_\alpha}{\xi^\beta} , \quad   
    \\ 
    b^\lambda_\mu &=  a^{\lambda \alpha} b_{\alpha \mu },\quad   g= \det (g_{\alpha \beta} ).
\end{align*}
In the following analysis we omit the dependence on the parameter $\bsy \in \calU$.
We now assume that the first two fundamental forms $a_{\alpha \beta}, b_{\alpha\beta}$ induced by $\phi$ are diagonal. This is the case, for example, for the shell of revolution mapping in Section \ref{sec:numerical-example}. 
Let 
\begin{align*}
\Phi&: \Omega^\epsilon \to \R^3, \vxi \mapsto \phi(\uxi) + \xi^3 \va_3(\uxi) , \\
h &: \Omega \to \Omega^\epsilon, \vec \eta \mapsto \begin{pmatrix}
    \eta_1 \\ \eta_2 \\ \epsilon \eta_3,
\end{pmatrix}\end{align*}
where $\vxi = (\uxi, \xi^3)^\top $. We note that 
\begin{align}\label{eq:eta3transform}
    \DD{u}{\xi^3} = \epsilon^{-1} \DD{u}{\eta_3}.
\end{align}
The prestress satisfies: there exist $\bar \sigma_{ij} : \Omega \to \mathbb{M}_{3}^2$ such that

 \begin{align}\label{eq:normalstressdecay1}
     \sigma_{ij}^* \vec a_{3i}  \vec a_{3j} &= \epsilon^2 \bar \sigma_{ij} \vec a_{3i} \vec a_{3j}, \\
     \label{eq:normalstressdecay2}
    \sigma_{ij}^* \vec a_{3i}  \vec a_{\alpha j} &= \epsilon \bar \sigma_{ij} \vec a_{3i} \vec a_{\alpha j}.
 \end{align} The normal component of the prestress decays to zero with thickness $\epsilon$. The decay of the normal stress components comes from the coefficient $\epsilon^{-1}$ 
 in the prestress source problem \eqref{eq:sourceproblem_weak_star} expansion $m_{-1}(\sigma,\bsv), m_{0}(\sigma,\bsv), m_{1}(\sigma,\bsv)$ as we see below.
 %and  expansion in $\epsilon$, as follows
 
Since we assume that $g_{\alpha\beta}$ is diagonal, we may simplify 
 \begin{align}\label{eq:covariant_vec_expa}
     \vec g^i = \frac 1 {g_{ii}}  \vec g _{i}
 \end{align}
 and expand 
 \begin{align} \label{eq:expa_gi}
     \vec g_i &= \vec a_i + \epsilon \vec g_{i}^{(1)} \\ 
      \vec g_{i}^{(1)}& = \begin{cases}
          \eta_3 b^{i}_i \vec a_i, &  i \in \{ 1,2\}\\
          0, & i =3 .
      \end{cases} 
 \end{align}
\paragraph{Exapnsion in $\epsilon$.} With \eqref{eq:covariant_vec_expa}, the bilinear forms can be written in terms of $\epsilon$, using \eqref{eq:eta3transform} and the expansion of the covariant basis \eqref{eq:expa_gi} to obtain
 \begin{align*}
     a(\bsu,\bsv) &= \int_{\Omega^\epsilon} C_{ijkl} ( \vec a_{r} + \epsilon \vec g^{(1)}_r )_j ( \vec a_{s} + \epsilon \vec g^{(1)}_s )_l  u_{i,r} v_{k,s} \frac{\sqrt g}{g_{rr} g_{ss}} \mathop{dV}, \\ 
     b(\bsu,\bsv) &= \int_{\Omega^\epsilon} \sigma^*_{ij} ( \vec a_{r} + \epsilon \vec g^{(1)}_r )_i ( \vec a_{s} + \epsilon \vec g^{(1)}_s )_j  u_{m,r} v_{m,s} \frac{\sqrt g}{g_{rr} g_{ss}} \mathop{dV},\\
     m(\sigma,\bsv) &= \int_{\Omega^\epsilon} \sigma_{ij} (\vec a_r + \epsilon \vec g_{r}^{(1)}) _j  \DD{v_i}{\xi^r} \sqrt{g} \mathop{dV}.
 \end{align*} 
  The next step is to transform \eqref{eq:eigenvaluebucklingproblem} to $\Omega$ and coordinates $\vec \eta$, and collect powers of $\epsilon$. 
  Then
  \begin{align*}
      a(\bsu,\bsv)  & = \epsilon^{-2} a_{-2}(\bsu,\bsv)  + \ldots + \epsilon^{2} a_{2}(\bsu,\bsv), \\ 
      b(\bsu,\bsv)  & =  b_{0}(\bsu,\bsv)  + \ldots + \epsilon^{2} b_{2}(\bsu,\bsv), \\ 
      m(\sigma,\bsv) &=\epsilon^{-1} m_{-1}(\sigma,\bsv)  + \ldots + \epsilon m_{1}(\sigma,\bsv) .
  \end{align*} 
    Note that the singularities in $b$ with respect to $\epsilon$ vanish due to decay in \eqref{eq:normalstressdecay1} of the  prestress in the normal direction. In order to see that the assumption on the prestress is justified, consider the $\epsilon$-expansion of the prestress bilinear form as  \begin{align*}
      m_{-1}(\sigma,\bsv) &=  \int_{\Omega} \sigma_{ij} \vec a_{3j} \DD{v_i}{\eta_3} \sqrt g \mathop{dV}, \\
     m_{0}(\sigma,\bsv) &=  \int_{\Omega} \sigma_{ij} \vec a_{\alpha j} \DD{v_i}{\eta_\alpha} \sqrt g \mathop{dV}, \\ 
      m_{1}(\sigma,\bsv) &=  \int_{\Omega} \sigma_{ij} \vec g_{\alpha j}^{(1)} \DD{v_i}{\eta_\alpha} \sqrt g \mathop{dV}.
  \end{align*}
  The singularity at $\epsilon \to 0$ forces the prestress to go in the normal direction of the surface to zero.
  More precisely, conditions \eqref{eq:normalstressdecay1}, \eqref{eq:normalstressdecay2} are satisfied.

The subsequent bilinear forms can be expanded and collected in a similiar fashion:
\begin{align*}
       a_{-2}(\bsu,\bsv) &= \int_{\Omega} C_{ijkl} \vec a_{3j}\vec a_{3l} \DD{u_i}{\eta_3} \DD{v_k}{\eta_3} \sqrt g \mathop{dV}, \\
         a_{-1}(\bsu,\bsv) &= \int_{\Omega} C_{ijkl} \left[ \vec a_{\alpha j}\vec a_{3l} \DD{u_i}{\eta_\alpha} \DD{v_k}{\eta_3} \frac1{g_{\alpha\alpha}} + \vec a_{3 j}\vec a_{\alpha l} \DD{u_i}{\eta_3} \DD{v_k}{\eta_\alpha} \frac1{g_{\alpha\alpha}} \right]\sqrt g \mathop{dV}, \\
           a_{0}(\bsu,\bsv) &= \int_{\Omega} C_{ijkl}  \left[ \vec a_{\alpha j}\vec a_{\beta l} \DD{u_i}{\eta_\alpha} \DD{v_k}{\eta_\beta } \frac1{g_{\alpha\alpha}g_{\beta\beta}} + \vec g^{(1)}_{\alpha j}\vec a_{3 l} \DD{u_i}{\eta_\alpha} \DD{v_k}{\eta_3} \frac1{g_{\alpha \alpha}} + \vec a_{3 j}\vec g^{(1)}_{\beta l} \DD{u_i}{\eta_3} \DD{v_k}{\eta_\beta } \frac1{g_{\beta\beta}}\right] \sqrt g \mathop{dV},\\
            a_{1}(\bsu,\bsv) &= \int_{\Omega} C_{ijkl} \left[ \vec g^{(1)}_{\alpha j}\vec a_{\beta l} \DD{u_i}{\eta_\alpha} \DD{v_k}{\eta_\beta}  + \vec a_{\alpha j}\vec g^{(1)}_{\beta l} \DD{u_i}{\eta_\alpha} \DD{v_k}{\eta_\beta}  \right] \frac{\sqrt g}{g_{\alpha\alpha}g_{\beta\beta}}\mathop{dV}, \\
            a_{2}(\bsu,\bsv) &= \int_{\Omega} C_{ijkl} \vec g^{(1)}_{\alpha j}\vec g^{(1)}_{\beta l} \DD{u_i}{\eta_\alpha} \DD{v_k}{\eta_\beta} \frac{\sqrt g}{g_{\alpha\alpha}g_{\beta\beta}}\mathop{dV},
    \end{align*}
    and the same expansion for $b(\bsu,\bsv)$ yields 
    \begin{align*}
        b_{0}(\bsu,\bsv) &=\int_{\Omega} \overline{\sigma}_{ij}\vec a_{ri}\vec a_{sj} \DD{u_m}{\eta_k} \DD{v_m}{\eta_s}\frac{\sqrt g}{g_{rr}g_{ss}}  \mathop{dV}   , \\ 
        b_{1}(\bsu,\bsv) &=\int_{\Omega} \overline{\sigma}_{ij}  \left[ \vec a_{r i}\vec g_{\alpha j}^{(1)} \DD{u_m}{\eta_r} \DD{v_m}{\eta_\alpha} +  \vec g_{\alpha i}^{(1)}\vec a_{r j} \DD{u_m}{\eta_\alpha} \DD{v_m}{\eta_r}\right] \frac{\sqrt g}{{g_{rr} g_{\alpha\alpha}} }\mathop{dV} , \\ 
        b_{2}(\bsu,\bsv) &=\int_{\Omega} \overline{\sigma}_{ij} \vec g_{\alpha i}^{(1)}\vec g_{\beta j}^{(1)} \DD{u_m}{\eta_\alpha} \DD{v_m}{\eta_\beta} \frac{\sqrt g}{g_{\alpha\alpha} g_{\beta\beta}} \mathop{dV} , \\ 
    \end{align*}
  Multiplying both sides of \eqref{eq:eigenvaluebucklingproblem} with $\epsilon^2$, and consider the leading terms
    \begin{align*}
       a_{-2}(\bsu,\bsv) + \ldots + \epsilon^4 a_2 &= \lambda( \epsilon^2 b_0(\bsu,\bsv) + \ldots + \epsilon^4 b_{2}(\bsu,\bsv)), \\
        a_{-2}(\bsu,\bsv) & \approx \lambda  \epsilon^2 b_0(\bsu,\bsv),
    \end{align*}
    
\printbibliography
%\bibliography{literature}

\end{document}